% !atlas-meta
% id: summa-contra-gentiles
% title: Summa Contra Gentiles (Of God and His Creatures), Book I
% author: aquinas
% authorName: Thomas Aquinas
% language: English
% sourceLanguage: Latin
% translator: Joseph Rickaby
% translationYear: 1905
% source: https://www.nd.edu/Departments/Maritain/etext/gc.htm
% license: Public Domain (translator Joseph Rickaby, S.J., d. 1929)
% status: partial
% !end-atlas-meta
% NOTE: This is Book I ("On God") of the four-book Summa Contra Gentiles, in Joseph
% Rickaby's public-domain translation "Of God and His Creatures" (1905), which is
% itself a somewhat abridged/annotated rendering. Books II-IV are NOT included.
% Source text retrieved from the Jacques Maritain Center (nd.edu) chapter files via
% the Internet Archive Wayback Machine. Chapters 17, 19, 27, 37, and 41 were not
% available in the archived source and are therefore absent.
\documentclass[12pt]{article}
\usepackage[utf8]{inputenc}
\usepackage{ebgaramond}
\usepackage[margin=1.2in]{geometry}
\usepackage{parskip}
\newcommand{\work}[1]{\begin{center}\LARGE\scshape #1\end{center}}
\newcommand{\attrib}[2]{\begin{center}\large #1\\[2pt]\small\itshape trans.\ #2\end{center}\vspace{1em}}
\newcommand{\para}[1][]{\par\noindent\ifx&#1&\else\textsuperscript{\footnotesize #1}~\fi}
\begin{document}
% !body-start
\work{Of God and His Creatures (Summa Contra Gentiles), Book I}

\attrib{Thomas Aquinas}{Joseph Rickaby (public domain)}

\section*{Chapter 1. The Function of the Wise Man}

\para[1] My mouth shall discuss truth, and my lips shall detest the ungodly (Prov. vii,7).

\para[2] ACCORDING to established popular usage, which the Philosopher considers should be our guide in the naming of things, they are called 'wise' who put things in their right order and control them well. Now, in all things that are to be controlled and put in order to an end, the measure of control and order must be taken from the end in view; and the proper end of everything is something good. Hence we see in the arts that art A governs and, as it were, lords it over art B, when the proper end of art B belongs to A. Thus the art of medicine lords it over the art of the apothecary, because health, the object of medicine, is the end of all drugs that the apothecary's art compounds. These arts that lord it over others are called 'master-building,' or 'masterful arts'; and the 'master-builders' who practise them arrogate to themselves the name of 'wise men.' But because these persons deal with the ends in view of certain particular things, without attaining to the general end of all things, they are called 'wise in this or that particular thing,' as it is said, 'As a wise architect I have laid the foundation' (I Cor. iii, 10); while the name of 'wise' without qualification is reserved for him alone who deals with the last end of the universe, which is also the first beginning of the order of the universe. Hence, according to the Philosopher, it is proper to the wise man to consider the highest causes.

\para[3] Now the last end of everything is that which is intended by the prime author or mover thereof. The prime author and mover of the universe is intelligence, as will be shown later (B. II, Chap. XXIII, XXIV). Therefore the last end of the universe must be the good of the intelligence, and that is truth. Truth then must be the final end of the whole universe; and about the consideration of that end wisdom must primarily be concerned. And therefore the Divine Wisdom, clothed in flesh, testifies that He came into the world for the manifestation of truth: For this was I born, and unto this I came into the World, to give testimony to the truth (John xvii, 37). The Philosopher also rules that the first philosophy is the science of truth, not of any and every truth, but of that truth which is the origin of all truth, and appertains to the first principle of the being of all things; hence its truth is the principle of all truth, for things are in truth as they are in being.

\para[4] It is one and the same function to embrace either of two contraries and to repel the other. Hence, as it is the function of the wise man to discuss truth, particularly of the first beginning, so it is his also to impugn the contrary error. Suitably therefore is the double function of the wise man displayed in the words above quoted from the Sapiential Book, namely, to study, and upon study to speak out the truth of God, which of all other is most properly called truth, and this is referred to in the words, My mouth shall discuss truth, and to impugn error contrary to truth, as referred to in the words, And my lips shall detest the ungodly.

\section*{Chapter 2. Of the Author's Purpose}

\para[5] OF all human pursuits, the pursuit of wisdom is the more perfect, the more sublime, the more useful, and the more agreeable. The more perfect, because in so far as a man gives himself up to the pursuit of wisdom, to that extent he enjoys already some portion of true happiness. Blessed is the man that shall dwell in wisdom (Ecclus xiv, 22). The more sublime, because thereby man comes closest to the likeness of God, who hath made all things in wisdom (Ps. ciii, 24). The more useful, because by this same wisdom we arrive at the realm of immortality. The desire of wisdom shall lead to an everlasting kingdom (Wisd. vi, 21). The more agreeable, because her conversation hath no bitterness, nor her company any weariness, but gladness and joy (Wisd. viii, 16).

\para[6] But on two accounts it is difficult to proceed against each particular error: first, because the sacrilegious utterances of our various erring opponents are not so well known to us as to enable us to find reasons, drawn from their own words, for the confutation of their errors: for such was the method of the ancient doctors in confuting the errors of the Gentiles, whose tenets they were readily able to know, having either been Gentiles themselves, or at least having lived among Gentiles and been instructed in their doctrines. Secondly, because some of them, as Mohammedans and Pagans, do not agree with us in recognising the authority of any scripture, available for their conviction, as we can argue against the Jews from the Old Testament, and against heretics from the New. But these receive neither: hence it is necessary to have recourse to natural reason, which all are obliged to assent to. But in the things of God natural reason is often at a loss.

\section*{Chapter 3. That the Truths which we confess concerning God fall under two Modes or Categories}

\para[7] BECAUSE not every truth admits of the same mode of manifestation, and "a well-educated man will expect exactness in every class of subject, according as the nature of the thing admits," as is very well remarked by the Philosopher (Eth. Nicom. I, 1094b), we must first show what mode of proof is possible for the truth that we have now before us. The truths that we confess concerning God fall under two modes. Some things true of God are beyond all the competence of human reason, as that God is Three and One. Other things there are to which even human reason can attain, as the existence and unity of God, which philosophers have proved to a demonstration under the guidance of the light of natural reason. That there are points of absolute intelligibility in God altogether beyond the compass of human reason, most manifestly appears. For since the leading principle of all knowledge of any given subject-matter is an understanding of the thing's innermost being, or substance --- according to the doctrine of the Philosopher, that the essence is the principle of demonstration --- it follows that the mode of our knowledge of the substance must be the mode of knowledge of whatever we know about the substance. Hence if the human understanding comprehends the substance of anything, as of a stone or triangle, none of the points of intelligibility about that thing will exceed the capacity of human reason. But this is not our case with regard to God. The human understanding cannot go so far of its natural power as to grasp His substance, since under the conditions of the present life the knowledge of our understanding commences with sense; and therefore objects beyond sense cannot be grasped by human understanding except so far as knowledge is gathered of them through the senses. But things of sense cannot lead our understanding to read in them the essence of the Divine Substance, inasmuch as they are effects inadequate to the power that caused them. Nevertheless our understanding is thereby led to some knowledge of God, namely, of His existence and of other attributes that must necessarily be attributed to the First Cause. There are, therefore, some points of intelligibility in God, accessible to human reason, and other points that altogether transcend the power of human reason.

\para[8] The same thing may be understood from consideration of degrees of intelligibility. Of two minds, one of which has a keener insight into truth than the other, the higher mind understands much that the other cannot grasp at all, as is clear in the 'plain man' (in rustico), who can in no way grasp the subtle theories of philosophy. Now the intellect of an angel excels that of a man more than the intellect of the ablest philosopher excels that of the plainest of plain men (rudissimi idiotae). The angel has a higher standpoint in creation than man as a basis of his knowledge of God, inasmuch as the substance of the angel, whereby he is led to know God by a process of natural knowledge, is nobler and more excellent than the things of sense, and even than the soul itself, whereby the human mind rises to the knowledge of God. But the Divine Mind exceeds the angelic much more than the angelic the human. For the Divine Mind of its own comprehensiveness covers the whole extent of its substance, and therefore perfectly understands its own essence, and knows all that is knowable about itself; but an angel of his natural knowledge does not know the essence of God, because the angel's own substance, whereby it is led to a knowledge of God, is an effect inadequate to the power of the cause that created it. Hence not all things that God understands in Himself can be grasped by the natural knowledge of an angel; nor is human reason competent to take in all that an angel understands of his own natural ability. As therefore it would be the height of madness in a 'plain man' to declare a philosopher's propositions false, because he could not understand them, so and much more would a man show exceeding folly if he suspected of falsehood a divine revelation given by the ministry of angels, on the mere ground that it was beyond the investigation of reason.

\para[9] The same thing manifestly appears from the incapacity which we daily experience in the observation of nature. We are ignorant of very many properties of the things of sense; and of the properties that our senses do apprehend, in most cases we cannot perfectly discover the reason. Much more is it beyond the competence of human reason to investigate all the points of intelligibility in that supreme excellent and transcendent substance of God. Consonant with this is the saying of the Philosopher, that "as the eyes of bats are to the light of the sun, so is the intelligence of our soul to the things most manifest by nature" (Aristotle, Metaphysics I, min. i). To this truth Holy Scripture also bears testimony. For it is said: Perchance thou wilt seize upon the traces of God, and fully discover the Almighty (Job xi, 7). And, Lo, God is great, and surpassing our knowledge (Job xxxvi, 26). And, We know in part (I Cor. xiii, 9). Not everything, therefore, that is said of God, even though it be beyond the power of reason to investigate, is at once to be rejected as false.

\section*{Chapter 4. That it is an advantage for the Truths of God, known by Natural Reason, to be proposed to men to be believed on faith}

\para[10] IF a truth of this nature were left to the sole enquiry of reason, three disadvantages would follow. One is that the knowledge of God would be confined to few. The discovery of truth is the fruit of studious enquiry. From this very many are hindered. Some are hindered by a constitutional unfitness, their natures being ill-disposed to the acquisition of knowledge. They could never arrive by study to the highest grade of human knowledge, which consists in the knowledge of God. Others are hindered by the needs of business and the ties of the management of property. There must be in human society some men devoted to temporal affairs. These could not possibly spend time enough in the learned lessons of speculative enquiry to arrive at the highest point of human enquiry, the knowledge of God. Some again are hindered by sloth. The knowledge of the truths that reason can investigate concerning God presupposes much previous knowledge. Indeed almost the entire study of philosophy is directed to the knowledge of God. Hence, of all parts of philosophy, that part stands over to be learnt last, which consists of metaphysics dealing with points of Divinity. Thus, only with great labour of study is it possible to arrive at the searching out of the aforesaid truth; and this labour few are willing to undergo for sheer love of knowledge. Another disadvantage is that such as did arrive at the knowledge or discovery of the aforesaid truth would take a long time over it, on account of the profundity of such truth, and the many prerequisites to the study, and also because in youth and early manhood, the soul, tossed to and fro on the waves of passion, is not fit for the study of such high truth: only in settled age does the soul become prudent and scientific, as the Philosopher says. Thus, if the only way open to the knowledge of God were the way of reason, the human race would dwell long in thick darkness of ignorance: as the knowledge of God, the best instrument for making men perfect and good, would accrue only to a few, and to those few after a considerable lapse of time.

\para[11] A third disadvantage is that, owing to the infirmity of our judgement and the perturbing force of imagination, there is some admixture of error in most of the investigations of human reason. This would be a reason to many for continuing to doubt even of the most accurate demonstrations, not perceiving the force of the demonstration, and seeing the divers judgements of divers persons who have the name of being wise men. Besides, in the midst of much demonstrated truth there is sometimes an element of error, not demonstrated but asserted on the strength of some plausible and sophistic reasoning that is taken for a demonstration. And therefore it was necessary for the real truth concerning divine things to be presented to men with fixed certainty by way of faith. Wholesome therefore is the arrangement of divine clemency, whereby things even that reason can investigate are commanded to be held on faith, so that all might easily be partakers of the knowledge of God, and that without doubt and error. Hence it is said: Now ye walk not as the Gentiles walk in the vanity of their own notions, having the understanding darkened (Eph. iv, 17, 18); and, I will make all thy sons taught of the Lord (Isa. liv, 1, 5).

\section*{Chapter 5. That it is an advantage for things that cannot he searched out by Reason to be proposed as Tenets of Faith}

\para[12] SOME may possibly think that points which reason is unable to investigate ought not to be proposed to man to believe, since Divine Wisdom provides for every being according to the measure of its nature; and therefore we must show the necessity of things even that transcend reason being proposed by God to man for his belief.

\para[13] 1. One proof is this. No one strives with any earnestness of desire after anything, unless it be known to him beforehand. Since, then, as will be traced out in the following pages (B.III, Chap.CXLVIII), Divine Providence directs men to a higher good than human frailty can experience in the present life, the mental faculties ought to be evoked and led onward to something higher than our reason can attain at present, learning thereby to desire something and earnestly to tend to something that transcends the entire state of the present life. And such is the special function of the Christian religion, which stands alone in its promise of spiritual and eternal goods, whereas the Old Law, carrying temporal promises, proposed few tenets that transcended the enquiry of human reason.

\para[14] 2. Also another advantage is thence derived, to wit, the repression of presumption, which is the mother of error. For there are some so presumptuous of their own genius as to think that they can measure with their understanding the whole nature of the Godhead, thinking all that to be true which seems true to them, and that to be false which does not seem true to them. In order then that the human mind might be delivered from this presumption, and attain to a modest style of enquiry after truth, it was necessary for certain things to be proposed to man from God that altogether exceeded his understanding.

\para[15] 3. There is also another evident advantage in this, that any knowledge, however imperfect, of the noblest objects confers a very high perfection on the soul. And therefore, though human reason cannot fully grasp truths above reason, nevertheless it is much perfected by holding such truths after some fashion at least by faith. And therefore it is said: Many things beyond the understanding of man are shown to thee (Ecclus iii, 23). And, The things that are of God, none knoweth but the Spirit of God: but to us God hath revealed them through his Spirit (I Cor. ii, 10, 11).

\section*{Chapter 6. That there is no lightmindedness in assenting to Truths of Faith, although they are above Reason}

\para[16] THE Divine Wisdom, that knows all things most fully, has deigned to reveal these her secrets to men, and in proof of them has displayed works beyond the competence of all natural powers, in the wonderful cure of diseases, in the raising of the dead, and what is more wonderful still, in such inspiration of human minds as that simple and ignorant persons, filled with the gift of the Holy Ghost, have gained in an instant the height of wisdom and eloquence. By force of the aforesaid proof, without violence of arms, without promise of pleasures, and, most wonderful thing of all, in the midst of the violence of persecutors, a countless multitude, not only of the uneducated but of the wisest men, flocked to the Christian faith, wherein doctrines are preached that transcend all human understanding, pleasures of sense are restrained, and a contempt is taught of all worldly possessions. That mortal minds should assent to such teaching is the greatest of miracles, and a manifest work of divine inspiration leading men to despise the visible and desire only invisible goods. Nor did this happen suddenly nor by chance, but by a divine disposition, as is manifest from the fact that God foretold by many oracles of His prophets that He intended to do this. The books of those prophets are still venerated amongst us, as bearing testimony to our faith. This argument is touched upon in the text: Which (salvation) having begun to be uttered by the Lord, was confirmed by them that heard him even unto us, God joining in the testimony by signs and portents and various distributions of the Holy Spirit (Heb. ii, 3, 4). This so wonderful conversion of the world to the Christian faith is so certain a sign of past miracles, that they need no further reiteration, since they appear evidently in their effects. It would be more wonderful than all other miracles, if without miraculous signs the world had been induced by simple and low-born men to believe truths so arduous, to do works so difficult, to hope for reward so high. And yet even in our times God ceases not through His saints to work miracles for the confirmation of the faith.

\section*{Chapter 7. That the Truth of reason is not contrary to the Truth of Christian Faith}

\para[17] THE natural dictates of reason must certainly be quite true: it is impossible to think of their being otherwise. Nor a gain is it permissible to believe that the tenets of faith are false, being so evidently confirmed by God. Since therefore falsehood alone is contrary to truth, it is impossible for the truth of faith to be contrary to principles known by natural reason.

\para[18] 2. Whatever is put into the disciple's mind by the teacher is contained in the knowledge of the teacher, unless the teacher is teaching dishonestly, which would be a wicked thing to say of God. But the knowledge of principles naturally known is put into us by God, seeing that God Himself is the author of our nature. Therefore these principles also are contained in the Divine Wisdom. Whatever therefore is contrary to these principles is contrary to Divine Wisdom, and cannot be of God.

\para[19] 3. Contrary reasons fetter our intellect fast, so that it cannot proceed to the knowledge of the truth. If therefore contrary informations were sent us by God, our intellect would be thereby hindered from knowledge of the truth: but such hindrance cannot be of God.

\para[20] 4. What is natural cannot be changed while nature remains. But contrary opinions cannot be in the same mind at the same time: therefore no opinion or belief is sent to man from God contrary to natural knowledge. And therefore the Apostle says: The word is near in thy heart and in thy mouth, that is, the word of faith which we preach (Rom. x, 8). But because it surpasses reason it is counted by some as contrary to reason, which cannot be. To the same effect is the authority of Augustine (Gen. ad litt. ii, 18) : " What truth reveals can nowise be contrary to the holy books either of the Old or of the New Testament." Hence the conclusion is evident, that any arguments alleged against the teachings of faith do not proceed logically from first principles of nature, principles of themselves known, and so do not amount to a demonstration; but are either probable reasons or sophistical; hence room is left for refuting them.

\section*{Chapter 8. Of the Relation of Human Reason to the first Truth of Faith *}

\para[21] THE things of sense, from whence human reason takes its beginning of knowledge, retain in themselves some trace of imitation of God, inasmuch as they are, and are good; yet so imperfect is this trace that it proves wholly insufficient to declare the substance of God Himself. Since every agent acts to the producing of its own likeness, effects in their several ways bear some likeness to their causes: nevertheless the effect does not always attain to the perfect likeness of the agent that produces it. In regard then to knowledge of the truth of faith, which can only be thoroughly known to those who behold the substance of God, human reason stands so conditioned as to be able to argue some true likenesses to it: which likenesses however are not sufficient for any sort of demonstrative or intuitive comprehension of the aforesaid truth. Still it is useful for the human mind to exercise itself in such reasonings, however feeble, provided there be no presumptuous hope of perfect comprehension or demonstration. With this view the authority of Hilary agrees, who says (De Trinitate, ii, 10), speaking of such truth : "In this belief start, run, persist; and though I know that you will not reach the goal, still I shall congratulate you as I see you making progress. But intrude not into that sanctuary, and plunge not into the mystery of infinite truth; entertain no presumptuous hope of comprehending the height of intelligence, but understand that it is incomprehensible."

\section*{Chapter 9. The Order and Mode of Procedure in this Work}

\para[22] THERE is then a twofold sort of truth in things divine for the wise man to study: one that can be attained by rational enquiry, another that transcends all the industry of reason. This truth of things divine I do not call twofold on the part of God, who is one simple Truth, but on the part of our knowledge, as our cognitive faculty has different aptitudes for the knowledge of divine things. To the declaration therefore of the first sort of truth we must proceed by demonstrative reasons that may serve to convince the adversary. But because such reasons are not forthcoming for truth of the second sort, our aim ought not to be to convince the adversary by reasons, but to refute his reasonings against the truth, which we may hope to do, since natural reason cannot be contrary to the truth of faith. The special mode of refutation to be employed against an opponent of this second sort of truth is by alleging the authority of Scripture confirmed from heaven by miracles. There are however some probable reasons available for the declaration of this truth, to the exercise and consolation of the faithful, but not to the convincing of opponents, because the mere insufficiency of such reasoning would rather confirm them in their error, they thinking that we assented to the truth of faith for reasons so weak.

\para[23] According then to the manner indicated we will bend our endeavour, first, to the manifestation of that truth which faith professes and reason searches out, alleging reasons demonstrative and probable, some of which we have gathered from the books of philosophers and saints, for the establishment of the truth and the confutation of the opponent. Then, to proceed from what is more to what is less manifest in our regard, we will pass to the manifestation of that truth which transcends reason, solving the arguments of opponents, and by probable reasons and authorities, so far as God shall enable us, declaring the truth of faith. Taking therefore the way of reason to the pursuit of truths that human reason can search out regarding God, the first consideration that meets us is of the attributes of God in Himself; secondly of the coming forth of creatures from God; thirdly of the order of creatures to God as to their last end.

\section*{Chapter 10. Of the Opinion of those who say that the Existence of God cannot he proved, being a Self-evident Truth}

\para[24] THIS opinion rests on the following grounds:

\para[25] 1. Those truths are self-evident which are recognised at once, as soon as the terms in which they are expressed are known. Such a truth is the assertion that God exists: for by the name 'God' we understand something greater than which nothing can be thought. This notion is formed in the understanding by whoever hears and understands the name 'God,' so that God must already exist at least in the mind. Now He cannot exist in the mind only: for what is in the mind and in reality is greater than that which is in the mind only; but nothing is greater than God, as the very meaning of the name shows: it follows that the existence of God is a self evident truth, being evidenced by the mere meaning of the name.

\para[26] 2. The existence of a being is conceivable, that could not be conceived not to exist; such a being is evidently greater than another that could be conceived not to exist. Thus then something greater than God is conceivable if He could be conceived not to exist; but anything 'greater than God' is against the meaning of the name 'God.' It remains then that the existence of God is a self-evident truth.

\para[27] 3. Those propositions are most self-evident which are either identities, as 'Man is man,' or in which the predicates are included in the definitions of the subjects, as 'Man is an animal.' But in God of all beings this is found true, that His existence is His essence, as will be shown later (Chap. XXII); and thus there is one and the same answer to the question 'What is He?' and 'Whether He is.' Thus then, when it is said 'God is,' the predicate is either the same with the subject or at least is included in the definition of the subject; and thus the existence of God will be a self-evident truth.

\para[28] 4. Things naturally known are self-evident: for the knowledge of them is not attained by enquiry and study. But the existence of God is naturally known, since the desire of man tends naturally to God as to his last end, as will be shown further on (B. 111, Chap. XXV).

\para[29] 5. That must be self-evident whereby all other things are known; but such is God; for as the light of the sun is the principle of all visual perception, so the divine light is the principle of all intellectual cognition.

\section*{Chapter 11. Rejection of the aforesaid Opinion, and Solution of the aforesaid Reasons}

\para[30] THE above opinion arises partly from custom, men being accustomed from the beginning to hear and invoke the name of God. Custom, especially that which is from the beginning, takes the place of nature; hence notions wherewith the mind is imbued from childhood are held as firmly as if they were naturally known and self-evident. Partly also it owes its origin to the neglect of a distinction between what is self-evident of itself absolutely and what is self-evident relatively to us. Absolutely indeed the existence of God is self-evident, since God's essence is His existence. But since we cannot mentally conceive God's essence, his existence is not self-evident relatively to us.

\para[31] 1. Nor is the existence of God necessarily self-evident as soon as the meaning of the name 'God' is known. First, because it is not evident, even to all who admit the existence of God, that God is something greater than which nothing can be conceived, since many of the ancients said that this world was God. Then granting that universal usage understands by the name 'God' something greater than which nothing can be conceived, it will not follow that there exists in rerum natura something greater than which nothing can be conceived. For 'thing' and "notion implied in the name of the thing" must answer to one another. From the conception in the mind of what is declared by this name 'God' it does not follow that God exists otherwise than in the mind. Hence there will be no necessity either of that something, greater than which nothing can be conceived, existing otherwise than in the mind; and from this it does not follow that there is anything in rerum natura greater than which nothing can be conceived. And so the supposition of the nonexistence of God goes untouched. For the possibility of our thought outrunning the greatness of any given object, whether of the actual or of the ideal order, has nothing in it to vex the soul of any one except of him alone who already grants the existence in rerum natura of something than which nothing can be conceived greater.

\para[32] 2. Nor is it necessary for something greater than God to be conceivable, if His non-existence is conceivable. For the possibility of conceiving Him not to exist does not arise from the imperfection or uncertainty of His Being, since His Being is of itself most manifest, but from the infirmity of our understanding, which cannot discern Him as He is of Himself, but only by the effects which He produces; and so it is brought by reasoning to the knowledge of Him.

\para[33] 3. As it is self-evident to us that the whole is greater than its part, so the existence of God is most self-evident to them that see the divine essence, inasmuch as His essence is His existence. But because we cannot see His essence, we are brought to the knowledge of His existence, not by what He is in Himself but by the effects which He works.

\para[34] 4. Man knows God naturally as he desires Him naturally. Now man desires Him naturally inasmuch as he naturally desires happiness, which is a certain likeness to the divine goodness. Thus it is not necessary that God, considered in Himself, should be naturally known to man, but a certain likeness of God. Hence man must be led to a knowledge of God through the likenesses of Him that are found in the effects which He works.

\para[35] 5. God is that wherein all things are known, not as though other things could not be known without His being known first, as happens in the case of self-evident principles, but because through His influence all knowledge is caused in us.

\section*{Chapter 12. Of the Opinion of those who say that the Existence of God is a Tenet of Faith alone and cannot he demonstrated}

\para[36] THE falseness of this opinion is shown to us as well by the art of demonstration, which teaches us to argue causes from effects, as also by the order of the sciences, for if there be no knowable substance above sensible substances, there will be no science above physical science; as also by the efforts of philosophers, directed to the proof of the existence of God; as also by apostolic truth asserting: The invisible things of God are clearly seen, being understood by the things that are made (Rom. i, 20). The axiom that in God essence and existence are the same is to be understood of the existence whereby God subsists in Himself, the manner of which is unknown to us, as also is His essence; not of the existence which signifies an affirmative judgement of the understanding. For in the form of such affirmative judgement the fact that there is a God falls under demonstration; as our mind is led by demonstrative reasons to form such a proposition declaratory of the existence of God. In the reasonings whereby the existence of God is demonstrated it is not necessary to assume for a premise the essence or quiddity of God: but instead of the quiddity the effect is taken for a premise, as is done in demonstrations a posteriori from effect to cause. All the names of God are imposed either on the principle of denying of God Himself certain effects of His power, or from some habitude of God towards those effects. Although God transcends sense and the objects of sense, nevertheless sensible effects are the basis of our demonstration of the existence of God. Thus the origin of our own knowledge is in sense, even of things that transcend sense.

\section*{Chapter 13. Reasons in Proof of the Existence of God}

\para[37] WE will put first the reasons by which Aristotle proceeds to prove the existence of God from the consideration of motion as follows.

\para[38] Everything that is in motion is put and kept in motion by some other thing. It is evident to sense that there are beings in motion. A thing is in motion because something else puts and keeps it in motion. That mover therefore either is itself in motion or not. If it is not in motion, our point is gained which we proposed to prove, namely, that we must posit something which moves other things without being itself in motion, and this we call God. But if the mover is itself in motion, then it is moved by some other mover. Either then we have to go on to infinity, or we must come to some mover which is motionless; but it is impossible to go on to infinity, therefore we must posit some motionless prime mover. In this argument there are two propositions to be proved: that everything which is in motion is put and kept in motion by something else; and that in the series of movers and things moved it is impossible to go on to infinity.

\para[39] The Philosopher also goes about in another way to show that it is impossible to proceed to infinity in the series of efficient causes, but we must come to one first cause, and this we call God. The way is more or less as follows. In every series of efficient causes, the first term is cause of the intermediate, and the intermediate is cause of the last. But if in efficient causes there is a process to infinity, none of the causes will be the first: therefore all the others will be taken away which are intermediate. But that is manifestly not the case; therefore we must posit the existence of some first efficient cause, which is God.

\para[40] Another argument is brought by St John Damascene (De Fid. Orthod. I, 3), thus: It is impossible for things contrary and discordant to fall into one harmonious order always or for the most part, except under some one guidance, assigning to each and all a tendency to a fixed end. But in the world we see things of different natures falling into harmonious order, not rarely and fortuitously, but always or for the most part. Therefore there must be some Power by whose providence the world is governed; and that we call God.

\section*{Chapter 14. That in order to a Knowledge of God we must use the Method of Negative Differentiation *}

\para[41] AFTER showing that there is a First Being, whom we call God, we must enquire into the conditions of His existence. We must use the method of negative differentiation, particularly in the consideration of the divine substance. For the divine substance, by its immensity, transcends every form that our intellect can realise; and thus we cannot apprehend it by knowing what it is, but we have some sort of knowledge of it by knowing what it is not. The more we can negatively differentiate it, or the more attributes we can strike off from it in our mind, the more we approach to a knowledge of it: for we know each thing more perfectly, the fuller view we have of its differences as compared with other things; for each thing has in itself a proper being, distinct from all others. Hence in dealing with things that we can define, we first place them in some genus, by which we know in general what the thing is; and afterwards we add the differentias whereby the thing is distinguished from other things; and thus is achieved a complete knowledge of the substance of the thing. But because in the study of the divine substance we cannot fix upon anything for a genus (Chap. ), nor can we mark that substance off from other things by affirmative differentias, we must determine it by negative differentias. In affirmative differentias one limits the extension of another, and brings us nearer to a complete designation of the thing under enquiry, inasmuch as it makes that thing differ from more and more things. And the same holds good also of negative differentias. For example, we may say that God is not an accident, in that He is distinguished from all accidents; then if we add that He is not a body, we shall further distinguish Him from some substances; and so in order by such negations He will be further distinguished from everything besides Himself; and then there will be a proper notion of His substance, when He shall be known as distinct from all. Still it will not be a perfect knowledge, because He will not be known for what He is in Himself.

\para[42] To proceed therefore in the knowledge of God by way of negative differentiation, let us take as a principle what has been shown in a previous chapter, that God is altogether immovable, which is confirmed also by the authority of Holy Scripture. For it is said: I am the Lord and change not (Mal. iii, 6) ; With whom there is no change (James i, 17); God is not as man, that he should change (Num. xxiii, 19).

\section*{Chapter 15. That God is Eternal}

\para[43] THE beginning of anything and its ceasing to be is brought about by motion or change. But it has been shown that God is altogether unchangeable: He is therefore eternal, without beginning or end.

\para[44] 2. Those things alone are measured by time which are in motion, inasmuch as time is an enumeration of motion. But God is altogether without motion, and therefore is not measured by time. Therefore in Him it is impossible to fix any before or after: He has no being after not being, nor can He have any not being after being, nor can any succession be found in His being, because all this is unintelligible without time. He is therefore without beginning and without end, having all His being at once, wherein consists the essence of eternity.

\para[45] 3. If at some time God was not, and afterwards was, He was brought forth by some cause from not being to being. But not by Himself, because what is not cannot do anything. But if by another, that other is prior to Him. But it has been shown that God is the First Cause; therefore He did not begin to be: hence neither will He cease to be; because what always has been has the force of being always.

\para[46] 4. We see in the world some things which are possible to be and not to be. But everything that is possible to be has a cause: for seeing that of itself it is open to two alternatives, being and not being; if being is to be assigned to it, that must be from some cause. But we cannot proceed to infinity in a series of causes: therefore we must posit something that necessarily is. Now everything necessary either has the cause of its necessity from elsewhere, or not from elsewhere, but is of itself necessary. But we cannot proceed to infinity in the enumeration of things necessary that have the cause of their necessity from elsewhere: therefore we must come to some first thing necessary, that is of itself necessary; and that is God. Therefore God is eternal, since everything that is of itself necessary is eternal.

\para[47] Hence the Psalmist: But thou, O Lord, abidest for ever: thou art the self-same, and thy years shall not fail (Ps. ci, 13-28).

\section*{Chapter 16. That in God there is no Passive Potentiality}

\para[48] EVERYTHING that has in its substance an admixture of potentiality, to the extent that it has potentiality is liable not to be: because what can be, can also not be. But God in Himself cannot not be, seeing that He is everlasting; therefore there is in God no potentiality.

\para[49] 2. Although in order of time that which is sometimes in potentiality, sometimes in actuality, is in potentiality before it is in actuality, yet, absolutely speaking, actuality is prior to potentiality, because potentiality does not bring itself into actuality, but is brought into actuality by something which is already in actuality. Everything therefore that is any way in potentiality has something else prior to it. But God is the First Being and the First Cause, and therefore has not in Himself any admixture of potentiality.

\para[50] 4. Everything acts inasmuch as it is in actuality. Whatever then is not all actuality, does not act by its whole self, but by something of itself. But what does not act by its whole self, is not a prime agent; for it acts by participation in something else, not by its own essence. The prime agent then, which is God, has no admixture of potentiality, but is pure actuality.

\para[51] 6. We see that there is that in the world which passes from potentiality to actuality. But it does not educe itself from potentiality to actuality, because what is in potentiality is not as yet, and therefore cannot act. Therefore there must be some other prior thing, whereby this thing may be brought out from potentiality to actuality. And again, if this further thing is going out from potentiality to actuality, there must be posited before it yet some other thing, whereby it may be reduced to actuality. But this process cannot go on for ever: therefore we must come to something that is only in actuality, and nowise in potentiality; and that we call God.

\section*{Chapter 18. That in God there is no Composition}

\para[52] IN every compound there must be actuality and potentiality. For a plurality of things cannot become one thing, unless there be actuality and potentiality. For things that are not one absolutely, are not actually united except by being in a manner tied up together or driven together: in which case the parts thus got together are in potentiality in respect of union; for they combine actually, after having been potentially combinable. But in God there is no potentiality: therefore there is not in Him any composition.

\para[53] 3. Every compound is potentially soluble in respect of its being compound, although in some cases there may be some other fact that stands in the way of dissolution. But what is soluble is in potentiality not to be, which cannot be said of God, seeing that He is of Himself a necessary Being.

\section*{Chapter 20. That God is Incorporeal}

\para[54] EVERY corporeal thing, being extended, is compound and has parts. But God is not compound: therefore He is not anything corporeal.

\para[55] 5. According to the order of objects is the order and distinction of powers: therefore above all sensible objects there is some intelligible object, existing in the nature of things. But every corporeal thing existing in nature is sensible: therefore there is determinable above all corporeal things something nobler than they. If therefore God is corporeal, He is not the first and greatest Being.

\para[56] With this demonstrated truth divine authority also agrees. For it is said: God is a spirit (John iv, 24): To the King of ages, immortal, invisible, only God (1 Tim. i, 17): The invisible things of God are understood and discerned by the things that are made (Rom. i, 29). For the things that are discerned, not by sight but by understanding, are incorporeal.

\para[57] Hereby is destroyed the error of the first natural philosophers, who posited none but material causes. The Gentiles also are refuted, who set up the elements of the world, and the powers therein existing, for gods; also the follies of the Anthropomorphite heretics, who figured God under bodily lineaments; also of the Manicheans, who thought God was an infinite substance of light diffused through infinite space. The occasion of all these errors was that, in thinking of divine things, men came under the influence of the imagination, which can be cognisant only of bodily likeness. And therefore we must transcend imagination in the study of things incorporeal.

\section*{Chapter 21. That God is His own Essence *}

\para[58] IN everything that is not its own essence, quiddity, or nature, there must be some composition. For since in everything its own essence is contained, --- if in anything there were contained nothing but its essence, the whole of that thing would be its essence, and so itself would be its own essence. If then anything is not its own essence, there must be something in that thing besides its essence, and so there must be in it composition. Hence also the essence in compound things is spoken of as a part, as humanity in man. But it has been shown that in God there is no composition. God therefore is His own essence.

\para[59] 2. That alone is reckoned to be beyond the essence of a thing, which does not enter into its definition: for the definition declares what the thing essentially is. But the accidents of a thing are the only points about it which fall not within the definition: therefore the accidents are the only points about a thing besides its essence. But in God there are no accidents, as will be shown (Chap. ): therefore there is nothing in Him besides His essence.

\para[60] 3. The forms that are not predicable of subsistent things, whether in the universal or in the singular, are forms that do not of themselves subsist singly, individualised in themselves. It is not said that Socrates or man or animal is whiteness; because whiteness is not anything subsisting singly in itself, but is individualised by the substance in which it exists. Also the essences or quiddities of genera or species are individualised according to the definite matter of this or that individual, although the generic or specific quiddity includes form and matter in general: hence it is not said that Socrates or man is humanity. But the Divine Essence is something existing singly by itself, and individualised in itself, as will be shown (Chap. XLII). The Divine Essence therefore is predicated of God in such a way that it can be said: 'God is His own essence.'

\section*{Chapter 22. That in God Existence and Essence is the same *}

\para[61] IT has been shown above (Chap. XV, n. 4) that there is an Existence which of itself necessarily is; and that is God. If this existence, which necessarily is, is contained in some essence not identical with it, then either it is dissonant and at variance with that essence, as subsistent existence is at variance with the essence of whiteness; or it is consonant with and akin to that essence, as existence in something other than itself is consonant with whiteness. In the former case, the existence which of itself necessarily is will not attach to that essence, any more than subsistent existence will attach to whiteness. In the latter case, either such existence must depend on the essence, or both existence and essence depend on another cause, or the essence must depend on the existence. The former two suppositions are against the idea of a being which of itself necessarily is; because, if it depends on another thing, it no longer is necessarily. From the third supposition it follows that that essence is accidental and adventitious to the thing which of itself necessarily is; because all that follows upon the being of a thing is accidental to it; and thus the supposed essence will not be the essence at all. God therefore has no essence that is not His existence.

\para[62] 2. Everything is by its own existence. Whatever then is not its own existence does not of itself necessarily exist. But God does of Himself necessarily exist: therefore God is His own existence.

\para[63] 4. 'Existence' denotes a certain actuality: for a thing is not said to 'be' for what it is potentially, but for what it is actually. But everything to which there attaches an actuality, existing as something different from it, stands to the same as potentiality to actuality. If then the divine essence is something else than its own existence, it follows that essence and existence in God stand to one another as potentiality and actuality. But it has been shown that in God there is nothing of potentiality (Chap. XVI), but that He is pure actuality. Therefore God's essence is not anything else but His existence.

\para[64] 5. Everything that cannot be except by the concurrence of several things is compound. But nothing in which essence is one thing, and existence another, can be except by the concurrence of several things, to wit, essence and existence. Therefore everything in which essence is one thing, and existence another, is compound. But God is not compound, as has been shown (Chap. XVIII). Therefore the very existence of God is His essence. This sublime truth was taught by the Lord to Moses (Exod. iii, 13, 14) If they say to me, What is his name? what shall I say to them? Thus shalt thou say to the children of Israel: He who is hath sent me to you: showing this to be His proper name, He who is. But every name is given to show the nature or essence of some thing. Hence it remains that the very existence or being of God is His essence or nature.

\section*{Chapter 23. That in God there is no Accident}

\para[65] EVERYTHING that is in a thing accidentally has a cause for its being therein, seeing that it is beside the essence of the thing wherein it is. If then there is anything in God accidentally, this must be by some cause. Either therefore the cause of the accident is the Divinity itself, or something else. If something else, that something must act upon the divine substance: for nothing induces any form, whether substantial or accidental, in any recipient, except by acting in some way upon it, because acting is nothing else than making something actually be, which is by a form. Thus God will be acted upon and moved by some agent, which is against the conclusions of Chapter XIII. But if the divine substance itself is the cause of the accident supposed to be in it, then, --- inasmuch as it cannot possibly be the cause of it in so far as it is the recipient of it, because at that rate the same thing in the same respect would actualise itself, --- then this accident, supposed to be in God, needs must be received by Him in one respect and caused by Him in another, even as things corporeal receive their proper accidents by the virtue of their matter, and cause them by their form. Thus then God will be compound, the contrary of which has been above proved.

\para[66] 4. In whatever thing anything is accidentally, that thing is in some way changeable in its nature: for accident as such may be and may not be in the thing in which it is. If then God has anything attaching to Him accidentally, it follows that He is changeable, the contrary of which has above been proved (Chap. , ).

\para[67] 5. A thing into which an accident enters, is not all and everything that is contained in itself: because accident is not of the essence of the subject. But God is whatever He has in Himself. Therefore in God there is no accident. --- The premises are proved thus. Everything is found more excellently in cause than in effect. But God is cause of all: therefore whatever is in Him is found there in the most excellent way possible. But what most perfectly attaches to a thing is the very thing itself. This unity of identity is more perfect than the substantial union of one element with another, e.g., of form with matter; and that union again is more perfect than the union that comes of one thing being accidentally in another. It remains therefore that God is whatever He has.

\para[68] Hence Augustine (De Trinitate, v, c. 4, n. 5): "There is nothing accidental in God, because there is nothing changeable or perishable." The showing forth of this truth is the confutation of sundry Saracen jurists, who suppose certain "ideas" superadded to the Divine Essence.

\section*{Chapter 24. That the Existence of God cannot he characterised by the addition of any Substantial Differentia *}

\para[69] IT is impossible for anything actually to be, unless all things exist whereby its substantial being is characterised. An animal cannot actually be without being either a rational or an irrational animal. Hence the Platonists, in positing Ideas, did not posit self-existent Ideas of genera, seeing that genera are characterised and brought to specific being by addition of essential differentias; but they posited self-existent Ideas of species alone, seeing that for the (further) characterising of species (in the individuals belonging to it) there is no need of essential differentias. If then the existence of God is characterised and receives an essential characteristic by the addition of something else, that existence will not of itself actually be except by having that other thing superadded to it. But the existence of God is His own very substance, as has been shown. It would follow that the substance of God could not actually be except by something supervening upon it; and thence the further conclusion would ensue that the substance of God is not of itself necessarily existent, the contrary of which has been shown above (Chap. , n. 4)

\para[70] 2. Everything that needs something superadded to enable it to be, is in potentiality in respect of that addition. Now the divine substance is not in any way in potentiality, as has been shown ), but God's own substance is God's own being. Therefore His existence cannot be characterised by any superadded substantial characteristic.

\section*{Chapter 25. That God is not in any Genus}

\para[71] EVERYTHING that is in any genus has something in it whereby the nature of the genus is characterised and reduced to species: for there is nothing in the genus that is not in some species of it. But this is impossible in God, as has been shown in the previous chapter.

\para[72] 2. If God is in any genus, He is either in the genus of accident or the genus of substance. He is not in the genus of accident, for an accident cannot be the first being and the first cause. Again, He cannot be in the genus of substance: for the substance that is a genus is not mere existence : otherwise every substance would be its own existence, since the idea of the genus is maintained in all that is contained under the genus: at that rate no substance would be caused by another, which is impossible (Chap. , ). But God is mere existence: therefore He is not in any genus.

\para[73] 3. Whatever is in a genus differs in point of existence from other things that are in the same genus: otherwise genus would not be predicated of several things. But all things that are in the same genus must agree in the quiddity, or essence, of the genus: because of them all genus is predicated so as to answer the question what (quid) each thing is. Therefore the existence of each thing that exists in a genus is something over and above the quiddity of the genus. But that is impossible in God.

\para[74] 4. Everything is placed in a genus by reason of its quiddity. But the quiddity of God is His own mere (full) existence . Now a thing is not ranked in a genus on the score of mere (bare) existence: otherwise 'being,' in the sense of mere (bare) existence, would be a genus. But that 'being' cannot be a genus is proved in this way. If 'being' were a genus, some differentia would have to be found to reduce it to species. But no differentia participates in its genus: I mean, genus is never comprehended in the idea of the differentia: because at that rate genus would be put twice over in the definition of the species. Differentia then must be something over and above what is understood in the idea of genus. Now nothing can be over and above what is understood by the idea of 'being'; since 'being' enters into the conceivability of all things whereof it is predicated, and thus can be limited by no differentia.

\para[75] Hence it is also apparent that God cannot be defined, because every definition is by genus and differentias. It is apparent also that there can be no demonstration of God except through some effect of His production: because the principle of demonstration is a definition of the thing defined.

\section*{Chapter 26. That God is not the formal or abstract being of all things}

\para[76] THINGS are not distinguished from one another in so far as they all have being, because in this they all agree. If therefore things do differ from one another, either 'being' itself must be specified by certain added differentias, so that different things have a different specific being; or things must differ in this that 'being' itself attaches to specifically different natures. The first alternative is impossible, because no addition can be made to 'being,' in the way that differentia is added to genus, as has been said (Chap. , n. 4). It remains therefore that things differ in that they have different natures, to which 'being' accrues differently. But the divine being is not something accessory to any nature, but is the very nature or essence of God (Chap. ). If therefore the divine being were the formal and abstract being of all things, all things would have to be absolutely one.

\para[77] 4. What is common to many is not anything over and above the many except in thought alone. For example, 'animal' is not anything over and above Socrates and Plato and other animals, except in the mind that apprehends the form of 'animal' despoiled of all individualising and specifying marks: for what is really animal is man: otherwise it would follow that in Plato there were several animals, to wit, animal in general, and man in general, and Plato himself. Much less then is bare being in general anything over and above all existing things, except in the mind alone. If then God be being in general, God will be nothing more than a logical entity, something that exists in the mind alone.

\para[78] This error is set aside by the teaching of Holy Scripture, which confesses God lofty and high (Isa. vi, 1), and that He is above all (Rom. ix, 5). For if He is the being of all, then He is something of all, not above all. The supporters of this error are also cast out by the same sentence which casts out idolaters, who gave the incommunicable name of God to stocks and stones (Wisd. xiv, 8, 21). For if God were the being of all, it would not be more truly said, 'A stone is a being,' than 'A stone is God.'

\para[79] What has led men into this error is a piece of faulty reasoning. For, seeing that what is common to many is specialised and individualised by addition, they reckoned that the divine being, to which no addition is made, was not any individual being, but was the general being of all things: failing to observe that what is common or universal cannot really exist without addition, but merely is viewed by the mind without addition. 'Animal' cannot be without 'rational' or 'irrational' as a differentia, although it may be thought of without these differentias. Moreover, though the universal be thought of without addition, yet not without susceptibility of addition. 'Animal' would not be a genus if no differentia could be added to it; and so of other generic names. But the divine being is without addition, not only in thought, but also in rerum natura; and not only without addition, but without even susceptibility of addition. Hence from this very fact, that He neither receives nor can receive addition, we may rather conclude that God is not being in general, but individual being: for by this very fact His being is distinguished from all other beings, that nothing can be added to it. (Chap. ).

\section*{Chapter 28. That God is Universal Perfection}

\para[80] AS all perfection and nobility is in a thing inasmuch as the thing is, so every defect is in a thing inasmuch as the thing in some manner is not. As then God has being in its totality, so not-being is totally removed from Him, because the measure in which a thing has being is the measure of its removal from not-being. Therefore all defect is absent from God: He is therefore universal perfection.

\para[81] 2. Everything imperfect must proceed from something perfect: therefore the First Being must be most perfect.

\para[82] 3. Everything is perfect inasmuch as it is in actuality; imperfect, inasmuch as it is in potentiality, with privation of actuality. That then which is nowise in potentiality, but is pure actuality, must be most perfect; and such is God.

\para[83] 4. Nothing acts except inasmuch as it is in actuality: action therefore follows the measure of actuality in the agent. It is impossible therefore for any effect that is brought into being by action to be of a nobler actuality than is the actuality of the agent. It is possible though for the actuality of the effect to be less perfect than the actuality of the acting cause, inasmuch as action may be weakened on the part of the object to which it is terminated, or upon which it is spent. Now in the category of efficient causation everything is reducible ultimately to one cause, which is God, of whom are all things. Everything therefore that actually is in any other thing must be found in God much more eminently than in the thing itself; God then is most perfect.

\para[84] Hence the answer given to Moses by the Lord, when he sought to see the divine face or glory: I will show thee all good (Exod. xxxiii, 19).

\section*{Chapter 29. How Likeness to God may be found in Creatures}

\para[85] EFFECTS disproportionate to their causes do not agree with them in name and essence. And yet some likeness must be found between such effects and their causes: for it is of the nature of an agent to do something like itself. Thus also God gives to creatures all their perfections; and thereby He has with all creatures a likeness, and an unlikeness at the same time. For this point of likeness, however, it is more proper to say that the creature is like God than that God is like the creature. For that is said to be like a thing, which possesses its quality or form. Since then that which is found to perfection in God is found in other beings by some manner of imperfect participation, the said point of likeness belongs to God absolutely, but not so to the creature. And thus the creature has what belongs to God, and is rightly said to be like to God: but it cannot be said that God has what belongs to the creature, nor is it fitting to say that God is like the creature; as we do not say that a man is like his picture, and yet his picture is rightly pronounced to be like him.

\section*{Chapter 30. What Names can be predicated of God}

\para[86] WE may further consider what may be said or not said of God, or what may be said of Him only, what again may be said of God and at the same time also of other beings. Inasmuch as every perfection of the creature may be found in God, although in another and a more excellent way, it follows that whatever names absolutely denote perfection without defect, are predicated of God and of other beings, as for instance, 'goodness,' 'wisdom,' 'being,' and the like. But whatever names denote such perfection with the addition of a mode proper to creatures, cannot be predicated of God except by way of similitude and metaphor, whereby the attributes of one thing are wont to be adapted to another, as when a man is called a 'block' for the denseness of his understanding. Of this sort are all names imposed to denote the species of a created thing, as 'man,' and 'stone': for to every species is due its own proper mode of perfection and being. In like manner also whatever names denote properties that are caused in things by their proper specific principles, cannot be predicated of God otherwise than metaphorically. But the names that express such perfections with that mode of supereminent excellence in which they appertain to God, are predicated of God alone, as for instance, 'Sovereign Good,' 'First Being,' and the like. I say that some of the aforesaid names imply perfection without defect, if we consider that which the name was imposed to signify. But if we consider the mode of signification, every name is attended with defect: for by a name we express things as we conceive them in our understanding: but our understanding, taking its beginning of knowledge from sensible objects, does not transcend that mode which it finds in such sensible objects. In them the form is one thing, and that which has the form another. The form, to be sure, is simple, but imperfect, as not subsisting by itself: while that which has the form subsists, but is not simple --- nay, is concrete and composite. Hence whatever our understanding marks as subsisting, it marks in the concrete: what it marks as simple, it marks, not as something that is, but as that whereby something is. And thus in every name that we utter, if we consider the mode of signification, there is found an imperfection that does not attach to God, although the thing signified may attach to God in some eminent way, as appears in the name 'goodness' and 'good.' 'Goodness' denotes something as not subsisting by itself: 'good,' something as concrete and composite. In this respect, then, no name befits God suitably except in respect of that which the name is imposed to signify. Such names therefore may be both affirmed and denied of God, affirmed on account of the meaning of the name, denied on account of the mode of signification. But the mode of supereminence, whereby the said perfections are found in God, cannot be signified by the names imposed by us, except either by negation, as when we call God 'eternal' or 'infinite,' or by reference or comparison of Him to other things, as when He is called the 'First Cause' or the 'Sovereign Good.' For we cannot take in (capere) of God what He is, but what He is not, and how other beings stand related to Him.

\section*{Chapter 31. That the Plurality of divine Names is not inconsistent with the Simplicity of the Divine Being predicated of God and of other Beings}

\para[87] THE perfections proper to other things in respect of their several forms must be attributed to God in respect of His productivity alone, which productivity is no other than His essence. Thus then God is called 'wise,' not only in respect of His producing wisdom, but because, in so far as we are wise, we imitate in some measure His productivity, which makes us wise. But He is not called 'stone,' though He has made stones, because in the name of 'stone' is understood a determinate mode of being wherein a stone is distinguished from God. Still a stone is an imitation of God its cause, in being, in goodness, and other such respects. Something of the sort may be found in the cognitive and active powers of man. The intellect by its one power knows all that the sentient part knows by several powers, and. much more besides. Also, the higher the intellect, the more it can know by one effort, to which knowledge an inferior intellect does not attain without many efforts. Again, the royal power extends to all those particulars to which the divers powers under it are directed.

\para[88] Thus also God by His one simple being possesses all manner of perfection, all that other beings compass by divers faculties --- yea, much more. Hereby the need is clear of many names predicated of God: for as we cannot know Him naturally otherwise than by arriving at Him from the effects which He produces, the names whereby we denote His perfections must be several and diverse, answering to the diverse perfections that are found in things. But if we could understand His essence as it is in itself, and adapt to it a name proper to it, we should express it by one name only, as is promised to those who shall behold Him in essence: In that day there shall be one Lord, and his name shall be one (Zach. xiv, 9).

\section*{Chapter 32. That nothing is predicated of God and other beings synonymously *}

\para[89] AN effect that does not receive a form specifically like the form whereby the agent acts, is incapable of receiving in synonymous predication the name taken from that form. But, of the things whereof God is cause, the forms do not attain to the species of the divine efficacy, since they receive piecemeal and in particular what is found in God simply and universally.

\para[90] 3. Everything that is predicated of several things synonymously, is either genus species, differentia, accidens, or proprium. But nothing is predicated of God as genus, as has been shown (Chap. ); and in like manner neither as differentia; nor again as species, which is made up of genus and differentia; nor can any accident attach to Him, as has been shown (Chap. ); and thus nothing is predicated of God either as accident or as proprium, for proprium is of the class of accidents. The result is that nothing is predicated synonymously of God and other beings.

\para[91] 6. Whatever is predicated of things so as to imply that one thing precedes and the other is consequent and dependent on the former, is certainly not predicated synonymously. Now nothing is predicated of God and of other beings as though they stood in the same rank, but it is implied that one precedes, and the other is consequent and dependent. Of God all predicates are predicated essentially. He is called 'being' to denote that He is essence itself; and 'good,' to denote that He is goodness itself. But of other beings predications are made to denote participation. Thus Socrates is called 'a man,' not that he is humanity itself, but one having humanity. It is impossible therefore for any predicate to be applied synonymously and in the same sense to God and other beings.

\section*{Chapter 33. That it is not at all true that the application of common Predicates to God and to Creatures involves nothing beyond a mere Identity of Name}

\para[92] WHERE there is a mere accidental identity of name, there is no order or respect implied of one thing to another, but quite by accident one name is applied to several different things. But this is not the case with the names applied to God and to creatures: for in such a community of names we have regard to the order of cause and effect (Chap. , ).

\para[93] 2. Moreover, there is some manner of likeness of creatures to God (Chap. ).

\para[94] 3. When there is no more than a mere identity of name between several things, we cannot be led from one of them to the knowledge of another; but from the attributes found in creatures we are led to a knowledge of the attributes of God (Chap. , ).

\para[95] 5. There is no use predicating any name of any thing unless by the name we come to understand something about the thing. But if names are predicated of God and creatures by a mere coincidence of sound, we understand by those names nothing whatever about God, seeing that the significations of those names are known to us only inasmuch as they apply to creatures: there would at that rate be no use in saying or proving of God that God is a good being, or anything else of the sort.

\para[96] If it is said that by such names we only know of God what He is not --- in that, e.g., He is called 'living' as not being of the genus of inanimate things --- at least it must be allowed that the predicate 'living,' applied to God and to creatures, agrees in the negation of the inanimate, and thus will be something more than a bare coincidence of name.

\section*{Chapter 34. That the things that are said God and Creatures are said analogously}

\para[97] THUS then from the foregoing arguments the conclusion remains that things said alike of God and of other beings are not said either in quite the same sense, or in a totally different sense, but in an analogous sense, that is, in point of order or regard to some one object. And this happens in two ways: in one way inasmuch as many things have regard to one particular, as in regard to the one point of health an animal is called 'healthy' as being the subject of health medicine is called 'healthful' as being productive of health; food is 'healthy,' being preservative of health; urine, as being a sign of health: in another way, inasmuch as we consider the order or regard of two things, not to any third thing, but to one of the two, as 'being' is predicated of substance and accident inasmuch as accident is referred to substance, not that substance and accident are referred to any third thing. Such names then as are predicated of God and of other beings are not predicated analogously in the former way of analogy --- for then we should have to posit something before God --- but in the latter way.

\para[98] In this matter of analogous predication we find sometimes the same order in point of name and in point of thing named, sometimes not the same. The order of naming follows the order of knowing, because the name is a sign of an intelligible concept. When then that which is prior in point of fact happens to be also prior in point of knowledge, there is one and the same priority alike in point of the concept answering to the name and of the nature of the thing named. Thus substance is prior to accident by nature, inasmuch as substance is the cause of accident; and prior also in knowledge, inasmuch as substance is put in the definition of accident; and therefore 'being' is predicated of substance before it is predicated of accident, alike in point of the nature of the thing and in point of the concept attaching to the name. But when what is prior in nature is posterior in knowledge, in such cases of analogy there is not the same order alike in point of the thing named and in point of the concept attaching to the name. Thus the power of healing, that is in healing remedies, is prior by nature to the health that is in the animal, as the cause is prior to the effect: but because this power is known from its effect, it is also named from its effect: hence, though 'healthful' or 'health- producing,' is prior in order of fast, yet the application of the predicate 'healthy' to the animal is prior in point of the concept attaching to the name. Thus then, because we arrive at the knowledge of God from the knowledge of other realities, the thing signified by the names that we apply in common to God and to those other realities --- the thing signified, I say, is by priority in God, in the mode proper to God: but the concept attaching to the name is posterior in its application to Him: hence He is said to be named from the effects which He causes.

\section*{Chapter 35. That the several Names predicated of God are not synonymous}

\para[99] THOUGH the names predicated of God signify the same thing, still they are not synonymous, because they do not signify the same point of view. For just as divers realities are by divers forms assimilated to the one simple reality, which is God, so our understanding by divers concepts is in some sort assimilated to Him, inasmuch as, by several different points of view, taken from the perfections of creatures, it is brought to the knowledge of Him. And therefore our understanding is not at fault in forming many concepts of one thing; because that simple divine being is such that things can be assimilated to it in many divers forms. According to these divers conceptions the understanding invents divers names, an assigns them to God-names which, though they denote one and the same thing, yet clearly are not synonymous, since they are not assigned from the same point of view. The same meaning does not attach to the name in all these cases, seeing that the name signifies the concept of the understanding before it signifies the thing understood.

\section*{Chapter 36. That the Propositions which our Understanding forms of God are not void of meaning}

\para[100] FOR all the absolute simplicity of God, not in vain does our understanding form propositions concerning Him, putting together and putting asunder. For though our understanding arrives by way of divers concepts to the knowledge of God, still it understands the absolute oneness of the object answering to all those concepts. Our mind does not attribute the manner of its understanding to the object is understood: thus it does not attribute immateriality to a stone, though it knows the stone immaterially. And therefore it asserts unity of the object by an affirmative proposition, which is a sign of identity, when it says, 'God is good': in which case any diversity that the composition shows is referable to the understanding, but unity to the thing understood. And on the same principle sometimes our mind forms a statement about God with some mark of diversity by inserting a preposition, as when it is said, 'Goodness is in God.' Herein is marked a diversity, proper to the understanding; and a unity, proper to the thing.

\section*{Chapter 38. That God is His own Goodness *}

\para[101] EVERY good thing, that is not its own goodness, is called good by participation. But what is called good by participation presupposes something else before itself, whence it has received the character of goodness. This process cannot go to infinity, as there is no processus in infinitum in a series of final causes: for the infinite is inconsistent with any end, while good bears the character of an end. We must therefore arrive at some first good thing, which is not good by participation in reference to anything else, but is good by its own essence; and that is God.

\para[102] 4. What is, may partake of something; but sheer being can partake of nothing. For that which partakes, is potentiality: but being is actuality. But God is sheer being, as has been proved (Chap. ) : He is not then good by participation, but essentially so.

\para[103] 5. Every simple being has its existence and what it is, in one: if the two were different, simplicity would be gone. But God is absolute simplicity, as has been shown (Chap. ): therefore the very goodness that is in Him is no other than His own very self.

\para[104] The same reasoning shows that no other good thing is its own goodness: wherefore it is said: None is good but God alone (Mark x, 18; Luke xviii, 19).

\section*{Chapter 39. That in God there can be no Evil}

\para[105] ESSENTIAL being, and essential goodness, and all other things that bear the name of 'essential,' contain no admixture of any foreign element; although a thing that is good may contain something else besides being and goodness, for there is nothing to prevent the subject of one perfection being the subject also of another. Everything is contained within the bounds of its essential idea in such sort as to render it incapable of containing within itself any foreign element. But God is goodness, not merely good. There cannot therefore be in Him anything that is not goodness, and so evil cannot be in Him at all.

\para[106] 3. As God is His own being, nothing can be said of God that signifies participation. If therefore evil could be predicated of Him, the predication would not signify participation, but essence. Now evil cannot be predicated of any being so as to be the essence of any: for to an essentially evil thing there would be wanting being, since being is good. There cannot be any extraneous admixture in evil, as such, any more than in goodness. Evil therefore cannot be predicated of God.

\para[107] 5. A thing is perfect in so far as it is in actuality: therefore it will be imperfect inasmuch as it is failing in actuality. Evil therefore is either a privation, or includes a privation, or is nothing. But the subject of privation is potentiality; and that cannot be in God: therefore neither can evil.

\para[108] This truth also Holy Scripture confirms, saying: God is light, and there is no darkness in Him, (I John i, 5) Far from God impiety, and iniquity from the Almighty (Job xxxiv, 10).

\section*{Chapter 40. That God is the Good of all Good}

\para[109] GOD in His goodness includes all goodnesses, and thus is the good of all good.

\para[110] 2. God is good by essence: all other beings by participation: therefore nothing can be called good except inasmuch as it bears some likeness to the divine goodness. He is therefore the good of all good. Hence it is said of the Divine Wisdom: There came to me all good things along with it (Wisd. vii, 11).

\section*{Chapter 42. That God is One}

\para[111] THERE cannot possibly be two sovereign goods. But God is the sovereign good. Therefore there is but one God.

\para[112] 2. God is all-perfect, wanting in no perfection. If then there are several gods, there must be several thus perfect beings. But that is impossible: for if to none of them is wanting any perfection, nor is there any admixture of imperfection in any, there will be nothing to distinguish them one from another.

\para[113] 7. If there are two beings, each necessarily existent, they must agree in point of necessary existence. Therefore they must be distinguished by some addition made to one only or to both of them; and thus either one or both must be composite. But no composite being exists necessarily of itself, as has been shown above (Chap. ). Therefore there cannot be several necessary beings, nor several gods.

\para[114] 9. If there are two gods, this name 'God' is predicated of each either in the same sense or in different senses. If in different senses, that does not touch the present question: for there is nothing to prevent anything from being called by any name in a sense different from that in which the name is ordinarily borne, if common parlance so allows. But if the predication is in the same sense, there must be in both a common nature, logically considered. Either then this nature has one existence in both, or it has two different existences. If it has one existence, they will be not two but one being: for there is not one existence of two beings that are substantially distinct. But if the nature has a different existence in each possessor, neither of the possessors will be his own essence, or his own existence, as is proper to God (Chap. ): therefore neither of them is that which we understand by the name of God.

\para[115] 12. If there are many gods, the nature of godhead cannot be numerically one in each. There must be therefore something to distinguish the divine nature in this and that god: but that is impossible, since the divine nature does not admit of addition or difference, whether in the way of points essential or of points accidental (Chap. , ).

\para[116] 13. Abstract being is one only: thus whiteness, if there were any whiteness in the abstract, would be one only. But God is abstract being itself, seeing that He is His own being (Chap. ). Therefore there can be only one God.

\para[117] This declaration of the divine unity we can also gather from Holy Writ. For it is said: Hear, O Israel, the Lord thy God is one Lord (Deut. vi, 4) And, One Lord, one faith (Eph. iv, 5).

\para[118] By this truth the Gentiles are set aside in their assertion of a multitude of gods. Yet it must be allowed that many of them proclaimed the existence of one supreme God, by whom all the other beings to whom they gave the name of gods had been created. They awarded the name of godhead to all everlasting substances, chiefly on the score of their wisdom and felicity and their government of the world. And this fashion of speech is found even in Holy Scripture, where the holy angels, or even men bearing the office of judges, are called gods: There is none like thee among gods, O Lord (Ps. lxxxv, 8); and, I have said, Ye are gods (Ps. lxxxi, 6). Hence the Manicheans seem to be in greater opposition to this truth in their maintenance of two first principles, the one not the cause of the other.

\section*{Chapter 43. That God is Infinite}

\para[119] INFINITY cannot be attributed to God on the score of multitude, seeing there is but one God. Nor on the score of quantitative extension, seeing He is incorporeal. It remains to consider whether infinity belongs to Him in point of spiritual greatness. Spiritual greatness may be either in power or in goodness (or completeness) of nature. Of these two greatnesses the one follows upon the other: for by the fact of a thing being in actuality it is capable of action. According then to the completeness of its actuality is the measure of the greatness of its power. Thus it follows that spiritual beings are called great according to the measure of their completeness, as Augustine says: "In things in which greatness goes not by bulk, being greater means being better" (De Trinit. vi, 9). But in God infinity can be understood negatively only, inasmuch as there is no term or limit to His perfection. And so infinity ought to be attributed to God.

\para[120] 2. Every actuality inhering in another takes limitation from that wherein it is: for what is in another is therein according to the measure of the recipient. An actuality therefore that is in none, is bounded by none: thus, if whiteness were self-existent, the perfection of whiteness in it would have no bounds till it attained all the perfection of whiteness that is attainable. But God is an actuality in no way existent in another: He is not a form inherent in matter; nor does His being inhere in any form or nature; since He is His own being, His own existence (Chap. ). The conclusion is that He is infinite.

\para[121] 4. Actuality is more perfect, the less admixture it has of potentiality. Every actuality, wherewith potentiality is blended, has bounds set to its perfection: while that which is without any blend of potentiality is without bounds to its perfection. But God is pure actuality without potentiality (Chap. ), and therefore infinite.

\para[122] 6. There cannot be conceived any mode in which any perfection can be had more perfectly than by him, who is perfect by his essence, and whose being is his own goodness. But such is God: therefore anything better or more perfect than God is inconceivable. He is therefore infinite in goodness.

\para[123] 7. Our intellect, in understanding anything, reaches out to infinity; a sign whereof is this, that, given any finite quantity, our intellect can think of something greater. But this direction of our intellect to the infinite would be in vain, if there were not something intelligible that is infinite. There must therefore be some infinite intelligible reality, which is necessarily the greatest of realities; and this we call God.

\para[124] 8. An effect cannot reach beyond its cause: now our understanding cannot come but of God, who is the First Cause. If then our understanding can conceive something greater than any finite being, the conclusion remains that God is not finite.

\para[125] 9. Every agent shows greater power in action, the further from actuality is the potentiality which it reduces to actuality, as there is need of greater power to warm water than to warm air. But that which is not at all, is infinitely distant from actuality, and is not in any way in potentiality: therefore if the world was made a fact from being previously no fact at all, the power of the Maker must be infinite. This argument avails to prove the infinity of the divine power even to the mind of those who assume the eternity of the world. For they acknowledge God to be the cause of the substantial being of the world, although they think that substance to have been from eternity, saying that the eternal God is the cause of an ever-existing world in the same way that a foot would be the cause of an everlasting foot-print, if it had been from eternity stamped on the dust. Still, even accepting the position thus defined, it follows that the power of God is infinite. For whether He produced things in time, according to us, or from eternity, according to them, there can be nothing in the world of reality that He has not produced, seeing that He is the universal principle of being; and thus He has brought things to be, without presupposition of any matter or potentiality. Now the measure of active power must be taken according to the measure of potentiality or passivity; for the greater the pre-existing or preconceived passivity, the greater the active power required to reduce it to complete actuality. The conclusion remains that, as finite power in producing an effect is conditioned on the potentiality of matter, the power of God, not being conditioned on any potentiality, is not finite, but infinite, and so is His essence infinite. To this truth Holy Scripture bears witness: Great is the Lord and exceedingly to he praised, and of his greatness there is no end (Ps. cxliv, 3).

\section*{Chapter 44. That God has Understanding}

\para[126] IN no order of causes is it found that an intelligent cause is the instrument of an unintelligent one. But all causes in the world stand to the prime mover, which is God, as instruments to the principal agent. Since then in the world there are found many intelligent causes, the prime mover cannot possibly cause unintelligently.

\para[127] 5. No perfection is wanting in God that is found in any kind of beings (Chap. ): nor does any manner of composition result in Him for all that (Chap. ). But among the perfections of creatures the highest is the possession of understanding: for by understanding a thing is in a manner all things, having in itself the perfections of all things.

\para[128] 6. Everything that tends definitely to an end, either fixes its own end, or has its end fixed for it by another: otherwise it would not tend rather to this end than to that. But the operations of nature tend to definite ends: the gains of nature are not made by chance: for if they were, they would not be the rule, but the exception, for chance is of exceptional cases. Since then physical agents do not fix their own end, because they have no idea of an end, they must have an end fixed for them by another, who is the author of nature. But He could not fix an end for nature, had He not Himself understanding.

\para[129] 7. Everything imperfect is derived from something perfect: for perfection is naturally prior to imperfection, as actuality to potentiality. But the forms that exist in particular things are imperfect, for the very reason that they do exist in particular, and not in the universality of their idea, or the fulness of their ideal being. They must therefore be derived from some perfect forms, which are not under particular limitations. Such forms cannot be other than objects of understanding, seeing that no form is found in its universality or ideal fulness, except in the understanding. Consequently such forms must be endowed with understanding, if they are to subsist by themselves: for only by that endowment can they be operative. God therefore, who is the first actuality existing by itself, whence all others are derived, must be endowed with understanding.

\para[130] This truth also is in the confession of Catholic faith : for it is said: He is wise of heart and mighty of power (Job ix, 4): With him is strength and wisdom (Ibid. xii, 16): Thy wisdom is made wonderful to me (Ps. cxxxviii, 6): O depth of riches, of wisdom and of knowledge of God (Rom. vi, 33).

\section*{Chapter 45. That in God the Understanding is His very Essence}

\para[131] TO understand is an act of an intelligent being, existing in that being, not passing out to anything external, as the act of warming passes out to the object warmed: for an intelligible object suffers nothing from being understood, but the intelligence that understands it is perfected thereby. But whatever is in God is the divine essence. Therefore the act of understanding in God is the divine essence.

\para[132] 5. Every substance is for the sake of its activity. If therefore the activity of God is anything else than the divine substance, His end will be something different from Himself; and thus God will not be His own goodness, seeing that the good of every being is its end.

\para[133] From the act of understanding in God being identical with His being, it follows necessarily that the act of His understanding is absolutely eternal and invariable, exists in actuality only, and has all the other attributes that have been proved of the divine being. God then is not potentially intelligent, nor does He begin anew to understand anything, nor does He undergo any change or composition in the process of understanding.

\section*{Chapter 46. That God understands by nothing else than by His own Essence}

\para[134] UNDERSTANDING is brought actually to understand by an impression made on the understanding, just as sense comes actually to feel by an impression made on sense. The impression made on the understanding then is to the understanding as actuality to potentiality. If therefore the divine understanding came to understand by any impression made on the understanding other than the understanding itself, the understanding would be in potentiality towards that impression, which, it has been shown, cannot be (Chap. , ).

\para[135] 3. Any impression on the understanding that is in the understanding over and above its essence, has an accidental being: by reason of which fact our knowledge reckons as an accident. But there can be no accident in God.

\para[136] Therefore there is not in His understanding any impression besides the divine essence itself.

\section*{Chapter 47. That God perfectly understands Himself}

\para[137] WHEN by an impression on the understanding that power is brought to bear on its object, the perfection of the intellectual act depends on two things: one is the perfect conformity of the impression with the thing understood: the other is the perfect fixing of the impression on the understanding: which perfection is the greater, the greater the power of the understanding to understand. Now the mere divine essence, which is the intelligible representation whereby the divine understanding understands, is absolutely one and the same with God Himself and with the understanding of God. God therefore knows Himself most perfectly.

\para[138] 6. The perfections of all creatures are found at their best in God. But of perfections found in creatures the greatest is to understand God: seeing that the intellectual nature is pre-eminent above other natures, and the perfection of intellect is the act of understanding, and the noblest object of understanding is God. God therefore understands Himself perfectly.

\para[139] This also is confirmed by divine authority, for the Apostle says: The spirit of God searcheth into even the deep things of God (i Cor. ii, 10).

\section*{Chapter 48. That God primarily and essentially knows Himself alone}

\para[140] THE Understanding is in potentiality in regard of its object, in so far as it is a different thing from that object. If therefore the primary and essential object of divine understanding be something different from God, it will follow that God is in potentiality in respect of some other thing, which is impossible (Chap. ).

\para[141] 5. A thing understood is the perfection of him who understands it: for an understanding is perfected by actually understanding, which means being made one with the object understood. If therefore anything else than God is the first object of His understanding, something else will be His perfection, and will be nobler than He, which is impossible.

\section*{Chapter 49. That God knows other things besides Himself}

\para[142] WE are said to know a thing when we know its cause. But God Himself by His essence is the cause of being to others. Since therefore He knows His own essence most fully, we must suppose that He knows also other beings.

\para[143] 3. Whoever knows anything perfectly, knows all that can be truly said of that thing, and all its natural attributes. But a natural attribute of God is to be cause of other things. Since then He perfectly knows Himself, He knows that He is a cause: which could not be unless He knew something also of what He has caused, which is something different from Himself, for nothing is its own cause.

\para[144] Gathering together these two conclusions, it appears that God knows Himself as the primary and essential object of His knowledge, and other things as seen in His essence.

\section*{Chapter 50. That God has a particular Knowledge of all things}

\para[145] EVERY agent that acts by understanding has a knowledge of what it does, reaching to the particular nature of the thing produced; because the knowledge of the maker determines the form of the thing made. But God is cause of things by His understanding, seeing that in Him to be and to understand are one. But everything acts inasmuch as it is in actuality. God therefore knows in particular, as distinct from other things, whatever He causes to be.

\para[146] 3. The collocation of things, distinct and separate, cannot be by chance, for it is in regular order. This collocation of things, then, distinct and separate from one another, must be due to the intention of some cause. It cannot be due to the intention of any cause that acts by physical necessity, because physical nature is determined to one line of acton. Thus of no agent, that acts by physical necessity, can the intention reach to many distinct effects, inasmuch as they are distinct. The distinct arrangement and collocation of things must proceed from the intention of some knowing cause. Indeed it seems the proper function of intellect to remark the distinction of things. It belongs therefore to the First Cause, which of itself is distinct from all others, to intend the distinct and separate collocation of all the materials of the Universe.

\para[147] 4. Whatever God knows, He knows most perfectly: for there is in Him all perfection (Chap. ). Now what is known only in general is not known perfectly: the main points of the thing are not known, the finishing touches of its perfection, whereby its proper being is completely realised and brought out. Such mere general knowledge is rather a perfectible than a perfect knowledge of a thing. If therefore God in knowing His essence knows all things in their universality, He must also have a particular knowledge of things.

\para[148] 8. Whoever knows any nature, knows whether that nature be communicable: for he would not know perfectly the nature of 'animal,' who did not know that it was communicable to many. But the divine nature is communicable by likeness. God therefore knows in how many ways anything may exist like unto His essence. Hence arises the diversity of types, inasmuch as they imitate in divers ways the divine essence. God therefore has a knowledge of things according to their several particular types. This also we are taught by the authority of canonical Scripture. God saw all things that he had made, and they were very good (Gen. i, 31). Nor is there any creature invisible in his sight, but all things are naked and open to his eyes (Heb. iv, 13).

\section*{Chapter 51. Some Discussion of the Question how there is in the Divine Understanding a Multitude of Objects}

\para[149] THIS Multitude cannot be taken to mean that many objects of understanding have a distinct being in God. For these objects of understanding would be either the same with the divine essence, and at that rate multitude would be posited in the essence of God, a doctrine above rejected on many grounds (Chap. ); or they would be additions made to the divine essence, and at that rate there would be in God some accident, which we have above shown to be an impossibility (Chap. ). Nor again can there be posited any separate existence of these intelligible forms, which seems to have been the position of Plato, who, by way of avoiding the above inconveniences, introduced the doctrine of Ideas. For the forms of physical things cannot exist without matter, as neither can they be understood without matter. And even supposing them so to exist, even this would not suffice to explain God understanding a multitude of objects. For, assuming the aforesaid forms to exist outside the essence of God, and that God could not understand the multitude of things without them, such understanding being requisite to the perfection of His intellect, it would follow that God's perfection in understanding depended on another being than Himself, and consequently His perfection in being, seeing that His being is His understanding: the contrary of all which has been shown (Chap. ). Moreover, assuming what shall be proved hereafter (Bk II, Chap. XV), that whatever is beyond the essence of God is caused by God, the above forms, if they are outside of God, must necessarily be caused by Him. But He is cause of things by His understanding, as shall be shown (Bk II, Chap. XXIII, XXIV). Therefore God's understanding of these intelligible forms is a natural prerequisite for the existence of such forms. God's understanding then of the multitude of creatures is not to be explained by the existence of many intelligible abstract forms outside of God.

\section*{Chapter 52. Reasons to show how the Multitude of intelligible Ideal Forms has no Existence except in the Divine Understanding}

\para[150] IT is not to be supposed that the multitude of intelligible ideal forms is in any other understanding save the divine, say, the understanding of an angel. For in that case the divine understanding would depend, at least for some portion of its activity, upon some secondary intellect, which is impossible: for as substances are of God, so also all that is in substances: hence for the being of any of these forms in any secondary intellect there is prerequired an act of the divine intelligence, whereby God is cause.

\para[151] 2. It is impossible for one intellect to perform an intellectual operation by virtue of another intellect being disposed to that operation: that intellect itself must operate, which is disposed so to do. The fact then of many intelligible forms being in some secondary intellect cannot account for the prime intellect knowing the multitude of such forms.

\section*{Chapter 53. How there is in God a Multitude of Objects of Understanding}

\para[152] AN external object, coming to be an object of our understanding, does not thereby exist in our understanding in its own proper nature: but the impression (species) of it must be in our understanding, and by that impression our understanding is actualised, or comes actually to understand. The understanding, actualised and 'informed' by such an impression, understands the 'thing in itself.' The act of understanding is immanent in the mind, and at the same time in relation with the thing understood, inasmuch as the aforesaid 'impression,' which is the starting-point of the intellectual activity, is a likeness of the thing understood. Thus informed by the impression (species) of the thing, the understanding in act goes on to form in itself what we may call an 'intellectual expression' (intentio) of the thing. This expression is the idea (ratio, ) of the thing, and so is denoted by the definition. So it must be, for the understanding understands alike the thing absent and the thing present; in which respect imagination and understanding agree. But the understanding has this advantage over the imagination, that it understands the thing apart from the individualising conditions without which the thing exists not in rerum natura. This could not be except for the understanding forming to itself the aforesaid 'expression.' This 'expression' (intentio) in the understanding, being, we may say, the term of the intellectual activity, is different from the 'intellectual impression' (species intelligibilis), which actualises the understanding and which must be considered the starting-point of intellectual activity; and yet both the one and the other, both the 'impression' (species) and the 'expression' (intentio), are likenesses of the 'thing in itself,' which is the object of the understanding. From the fact of the intellectual impression, which is the form of the intellect and the starting-point of intellectual knowledge, being a likeness of the external thing, it follows that the expression, or idea, formed by the understanding, is also like the thing: for as an agent is, so are its activities. And again, from the fact of the expression, or idea, in the understanding being like to its object, it follows that the understanding in the act of forming such an idea understands the said object.

\para[153] But the divine mind understands by virtue of no impression other than its own essence (Chap. ). At the same time the divine essence is the likeness of all things. It follows therefore that the concept of the divine understanding itself, which is the Divine Word, is at once a likeness of God Himself understood, and also a likeness of all things whereof the divine essence is a likeness. Thus then by one intelligible impression (species intelligibilis), which the divine essence, and by one intellectual recognition (intentio intellecta), which is the Divine Word, many several objects may be understood by God.

\section*{Chapter 54. }

\para[154] That the Divine Essence, being One, is the proper Likeness and Type of all things Intelligible

\para[155] BUT again it may seem to some difficult or impossible that one and the same simple being, as the divine essence, should be the proper type (propria ratio) and likeness of different things. For as different things are distinguished by means of their proper forms, it needs must be that what is like one thing according to its proper form should be found unlike to another.

\para[156] True indeed, different things may have one point of likeness in so far as they have one common feature, as man and ass, inasmuch as they are animals. If it were by mere discernment of common features that God knew things, it would follow that He had not a particular but only a general knowledge of things (contrary to Chap. ). To return then to a proper and particular knowledge, of which there is here question. The act of knowledge is according to the mode in which the likeness of the known object is in the knowing mind: for the likeness of the known object in the knowing mind is as the form by which that mind is set to act. If therefore God has a proper and particular knowledge of many different things, He must be the proper and particular type of each. We have to enquire how that can be.

\para[157] As the Philosopher says, the forms of things, and the definitions which mark such forms, are like numbers, in which the addition or subtraction of unity varies the species of the number. So in definitions: one differentia subtracted or added varies the species: thus 'sentient substance' varies in species by the addition of 'irrational' or 'rational.' But in instances of 'the many in one' the condition of the understanding is not as the condition of concrete nature. The nature of a concrete being does not admit of the severance of elements, the union of which is requisite to the existence of that being: thus animal nature will not endure if the soul be removed from the body. But the understanding can sometimes take separately elements that in actual being are united, when one of them does not enter into the concept of the other; thus in 'three' it may consider 'two' only, and in 'rational animal' the 'sentient' element alone. Hence the understanding may take what is inclusive of many elements for a proper specimen of many, by apprehending some of them without others. It may take 'ten' as a proper specimen of nine by subtraction of one unit, and absolutely as a proper specimen of all the numbers included in 'ten.' So also in 'man' it might recognise a proper type of 'irrational animal' as such, and of all the species of 'irrational animal,' unless these species involved some positive differentias. Therefore a certain philosopher, named Clement, said that in the scale of beings the nobler are types and patterns of the less noble. Now the divine essence contains in itself the noble qualities of all beings, not by way of a compound but by way of a perfect being (Chap. ). Every form, as well particular as general, is a perfection in so far as it posits something; and involves imperfection only in so far as it falls short of true being. The divine understanding then can comprehend whatever is proper to each in its essence, by understanding wherein each thing imitates the divine essence, and wherein it falls short of the perfection proper to that essence. Thus, by understanding its own essence as imitable in the way of life without consciousness, it gathers the proper form of a plant, by understanding the same essence as imitable in the way of consciousness without intellect, the proper form of an animal; and so of the rest. Evidently then the divine essence, inasmuch as it is absolutely perfect, may be taken as the proper type of each entity; and hence by it God may have a particular knowledge of all. But because the proper type of one is distinct from the proper type of another --- and distinction is the principle of plurality --- there must be observable in the divine intellect a distinction and plurality of recognised types, in so far as the content of the divine mind is the proper type of different things. And as it is in this way that God is cognisant of the special relation of likeness that each creature bears to Him, it follows that the types (rationes) of things on the divine mind are not several or distinct, except in so far as God knows things to be in several divers ways capable of assimilation to Himself.

\para[158] And from this point of view Augustine says that God has made man in one plan and horse on another; and that the plans or types of things exist severally in the divine mind (De div. quaest., LXXXIII, 46). And herein also is defensible in some sort the opinion of Plato, who supposes Ideas, according to which all beings in the material world are formed.

\section*{Chapter 55. That God understands all things at once and together}

\para[159] THE reason why our understanding cannot understand many things together in one act is because in the act of understanding the mind becomes one with the object understood; whence it follows that, were the mind to understand many things together in one act, it would be many things together, all of one genus, which is impossible. Intellectual impressions are all of one genus: they are of one type of being in the existence which they have in the mind, although the things of which they are impressions do not agree in one type of being: hence the contrariety of things outside the mind does not render the impressions of those things in the mind contrary to one another. And hence it is that when many things are taken together, being anyhow united, they are understood together. Thus a continuous whole is understood at once, not part by part; and a proposition is understood at once, not first the subject and then the predicate: because all the parts are known by one mental impression of the whole. Hence we gather that whatever several objects are known by one mental presentation, can be understood together: but God knows all things by that one presentation of them, which is His essence; therefore He can understand all together and at once.

\para[160] 2. The faculty of knowledge does not know anything actually without some attention and advertence. Hence the phantasms, stored in the sensorium, are at times not actually in the imagination, because no attention is given to them. We do not discern together a multitude of things to which we do not attend together: but things that necessarily fall under one and the same advertence and attention, are necessarily understood together. Thus whoever institutes a comparison of two things, directs his attention to both and discerns both together. But all things that are in the divine knowledge must necessarily fall under one advertence; for God is attentive to behold His essence perfectly, which is to see it to the whole reach of its virtual content, which includes all things. God therefore, in beholding His essence, discerns at once all things that are.

\para[161] 6. Every mind that understands one thing after another, is sometimes potentially intelligent, sometimes actually so; for while it understands the first thing actually, it understands the second potentially. But the divine mind is never potentially intelligent, but always actually: it does not, then, understand things in succession, but all at once.

\para[162] Holy Scripture witnesses to this truth, saying that with God there is no change nor shadow of vicissitude (James i, 17).

\section*{Chapter 56. That there is no Habitual Knowledge in God}

\para[163] IN whatever minds there is habitual knowledge, not all things are known together: but some things are known actually, others habitually. But in God all things are known actually (Chap. ).

\para[164] 2. He who has a habit of knowledge, and is not adverting to what he knows, is in a manner in potentiality, although otherwise than as he was before he understood at all: but the divine mind is nowise in potentiality.

\para[165] 3. In every mind that knows anything habitually, the mind's essence is different from its intellectual activity, which is the act of attentive thought. To such a mind, in habitual knowledge, activity is lacking, though the essence of the mind itself cannot be lacking. But in God His essence is His activity (Chap. ).

\para[166] 4. A mind that knows habitually only, is not in its ultimate perfection: hence that best of goods, happiness, is not taken to be in habit but in act. If then God is habitually knowing, He will not be all-perfect (contrary to Chap. ).

\para[167] 5. As shown in chapter , God has understanding by His essence, not by any intelligible forms superadded to His essence. But every mind in habitual knowledge understands by some such forms: for a habit is either a predisposition of the mind to receive mental impressions, or forms, whereby it comes actually to understand; or it is an orderly aggregation of such forms, existing in the mind, not in complete actuality, but in some manner intermediate between potentiality and actuality.

\para[168] 6. A habit is a quality: but in God there can be neither quality nor any other accident (Chap. ): habitual knowledge therefore is not proper to God. Because the mental state of thinking, or willing, or acting habitually only, is like the state of a sleeper, David says, by way of removing all habitual states from God: Lo, he shall not slumber or sleep who keepeth Israel (Ps. cxx, 4). And again it is said: The eyes of the Lord are far brighter than the sun (Ecclus xxiii, 28), for the sun is always in the act of shining.

\section*{Chapter 57. That the Knowledge of God is not a Reasoned Knowledge}

\para[169] OUR thought is then reasoned, when we pass from one object of thought to another, as in making syllogisms from principles to conclusions. Reasoning or arguing does not consist in seeing how a conclusion follows from premises by inspection of both together. That is not argument, but judging of argument. Now God does not think of one thing after another in any sort of succession, but of all things at once (Chap. ). His knowledge therefore is not reasoned or argumentative, although He knows the argument and reason of all things.

\para[170] 2. Every reasoner intues principles with one thought, and the conclusion with another. There would be no need to proceed to a conclusion from the consideration of premises, if the mere consideration of the premises at once laid the conclusion bare. But God knows all things by one act which is His essence (Chap. ). His knowledge therefore is not argumentative.

\para[171] 3. All argumentative knowledge has something of actuality and something of potentiality, for conclusions are potentially in premises. But in the divine mind potentiality has no place.

\para[172] 5. Things that are known naturally are known without reasoning, as appears in the case of first principles. But in God there can be no knowledge that is not natural, nay, essential: for His knowledge is His essence.

\para[173] 7. Only in its highest advance does the inferior touch upon the superior. But the highest advance of our knowledge is not reasoning, but intuition (intellectus), which is the starting-point of reasoning. God's knowledge then is not 'rational,' in the sense of 'argumentative,' but intuitive only.

\para[174] 8. Reasoning means a lack of intuition: the divine knowledge therefore is not reasoned.

\para[175] If any should take it amiss that God cannot make a syllogism, let them mark that He has the knowledge how to make syllogisms as one judging of them, not as one arguing syllogistically.

\para[176] To this there is witness of Holy Scripture in the text: All things are naked and open to his eyes (Heb. iv, i 3): whereas things that we know by reasoning are not of themselves naked and open to us, but are opened out and laid bare by reason.

\section*{Chapter 58. That God does not understand by Combination and Separation of Ideas}

\para[177] THINGS mentally combinable and separable are naturally considered by the mind apart from one another: for there would be no need of their combination and separation, if by the mere apprehension of a thing as being it were at once understood what was in it or not in it. If therefore God understood by a mental process of combination and separation, it would follow that He did not take in all things at one glance, but each thing apart, contrary to what has been shown above (Chap. ).

\para[178] 3. A mind that combines and separates, forms different judgements by different combinations. For a mental combination does not go beyond the terms of the combination. Hence, in the combination, or affirmative judgement (compositione), whereby the mind judges that man is an animal, it does not judge that a triangle is a figure. Now combination or separation is an operation of the mind. If God therefore views things by mentally combining and separating them, His mental act will not be one only but manifold; and so His essence will not be one only.

\para[179] Not for this however must we say that He is ignorant of tenable propositions: for His one and simple essence is the pattern of all things manifold and compound; and so by it God knows the whole multitude and complexity as well of actual nature as of the ideal world (tam naturae quam rationis).

\para[180] This is in consonance with the authority of Holy Scripture: for it is said, For my thoughts are not your thoughts (Isa. lv, 8); and yet, The Lord knoweth the thoughts of men (Ps. xciii, 11), which certainly proceed by combination and separation of ideas.

\section*{Chapter 59. That the Truth to be found in Propositions is not excluded from God}

\para[181] ALTHOUGH the knowledge of the divine mind is not after the manner of combination and separation of ideas in affirmative and negative propositions, nevertheless there is not excluded from it that truth which, according to the Philosopher, obtains only in such combinations and separations. For since the truth of the intellect is an equation of the intellect and the thing, inasmuch as the intellect says that to be which is, or that not to be which is not, truth belongs to that in the intellect which the intellect says, not to the act whereby it says it; for it is not requisite to the truth of the intellect that the mere act of understanding be equated to the thing, but what the mind says and knows by understanding must be equated to the thing, so that the case of the thing shall be as the mind says it is. But God by his simple understanding, in which there is no combination and separation of ideas, knows not only the essence of things, but also the propositions that are tenable concerning them (Chap. , ). Thus what the divine mind says by understanding is affirmation and negation. Therefore the simplicity of the divine mind does not import the shutting out from it of truth.

\section*{Chapter 60. That God is Truth}

\para[182] TRUTH is a perfection of the understanding and of its act. But the understanding of God is His substance; and the very act of understanding, as it is the being of God, is perfect as the being of God is perfect, not by any superadded perfection, but by itself. It remains therefore that the divine substance is truth itself.

\para[183] 4. Though truth is properly not in things but in the mind, nevertheless a thing is sometimes called true, inasmuch as it properly attains the actuality of its proper nature. Hence Avicenna says that the truth of a thing is a property of the fixed and appointed being of each thing, inasmuch as such a thing is naturally apt to create a true impression of itself, and inasmuch as it expresses the proper idea of itself in the divine mind. But God is His own essence: therefore, whether we speak of truth of the intellect or truth of the object, God is His own truth.

\para[184] This is also confirmed by the authority of our Lord saying of Himself: I am the way and the truth and the life (John xiv, 6).

\section*{Chapter 61. That God is pure Truth}

\para[185] THE understanding is not liable to error in its knowledge of abstract being, as neither is sense in dealing with the proper object of each sense. But all the knowledge of the divine mind is after the manner of a mind knowing abstract being (Chap. ): it is impossible therefore for error or deception or falsehood to creep into the cognitive act of God.

\para[186] 3. The intellect does not err over first principles, but over reasoned conclusions from first principles. But the divine intellect is not reasoning or argumentative (Chap. ), and is therefore not liable to deception.

\para[187] 4. The higher any cognitive faculty is, the more universal and far-reaching is its proper object: hence what sight is cognisant of accidentally, general sensibility or imagination seizes upon as a content of its proper object. But the power of the divine mind is the acme of cognitive power: therefore all things knowable stand to it as proper and ordinary objects of knowledge, not as accidental objects. But over proper and ordinary objects of knowledge a cognitive faculty never makes a mistake.

\para[188] 5. An intellectual virtue is a perfection of the understanding in knowing. It never happens that the understanding utters anything false, but its utterance is always true, when prompted by any intellectual virtue; for it is the part of virtue to render an act good, and to utter truth is the good act of the understanding. But the divine mind, being the acme of perfection, is more perfect by its nature than the human mind by any habit of virtue.

\para[189] 6. The knowledge of the human mind is in a manner caused by things: hence it comes to be that things knowable are the measure of human knowledge: for the judgement of the mind is true, because the thing is so. But the divine mind by its knowledge is the cause of things. Hence God's knowledge must be the measure of things, as art is the measure of products of art, whereof the perfection of each varies according to its agreement with art. Thus the divine mind stands to things as things stand to the human mind. But any error that arises out of any inequality between the human mind and the thing is not in things, but in the mind. If therefore there were not an absolutely perfect correspondence of the divine mind with things, the error would be in the things, not in the divine mind. There is however no error in the things that be: because each has so much of truth as it has of being. There is then no failure of correspondence between the divine mind and the things that be.

\para[190] Hence it is said: God is truthful (Rom. iii, 4): God is not like man, that he should lie (Num. xxiii, 19): God is light, and there is no darkness in him (1 John i, 5).

\section*{Chapter 62. That the Truth of God is the First and Sovereign Truth}

\para[191] THE standard in every genus is the most perfect instance of the genus. But the divine truth is the standard of all truth. The truth of our mind is measured by the object outside the mind: our understanding is called true, inasmuch as it is in accordance with that object. And again the truth of the object is measured by its accordance with the divine mind, which is the cause of all things (B. II, Chap. XXIV), as the truth of artificial objects is measured by the art of the artificer. Since then God is the first understanding and the first object of understanding, the truth of every understanding must be measured by His truth, as everything is measured by the first and best of its kind.

\section*{Chapter 63. The Arguments of those who wish to withdraw from God the Knowledge of Individual Things *}

\para[192] THE first argument is drawn from the very condition of individuality. For as matter (materia signata) is the principle of individuality, it seems that individuals cannot be known by any immaterial faculty, inasmuch as all knowledge is a certain assimilation, and hence even in us those powers alone apprehend individual objects, that make use of material organs, as do the imagination and senses, but our understanding, which is immaterial, does not recognise individuals as such: much less then is the divine understanding apt to take cognisance of individuals, being, as it is, the furthest removed from matter.

\para[193] 2. The second argument is that individual things do not always exist. Either then they will always be known by God, or they will sometimes be known and sometimes not known. The former alternative is impossible, because there can be no knowledge of that which is not: for knowledge is only of things true, and things that are not cannot be true. The second alternative is also impossible, because the knowledge of the divine mind is absolutely invariable (Chap. ).

\para[194] 3. The third argument is from the consideration that not all individual things come of necessity, but some are by contingency: hence there can be no certain knowledge of them except when they exist. For that knowledge is certain, which is infallible: but all knowledge of contingent being is fallible while the thing is still in the future; for the opposite may happen of that which is held in cognition: for if the opposite could not happen, the thing would be a necessity: hence there can be no science in us of future contingencies, only a conjectural reckoning. On the other hand we must suppose that all God's knowledge is most certain and infallible (Chap. ). It is also impossible for God to begin to know anything, by reason of His immutability. From this it seems to follow that He does not know individual contingencies.

\para[195] 4. The fourth argument is from this, that some individual effects have their cause in will. Now an effect, before it is produced, can be known only in its cause: for so only can it have being before it begins to have being in itself. But the motions of the will can be known with certainty by none other than the willing agent, in whose power they are. It is impossible therefore that God should have certain knowledge of such individual effects as derive their causation from a created will.

\para[196] 5. The fifth argument is from the infinite multitude of individual things. The infinite as such is unknown: for all that is known is measured in a manner by the comprehension of the knower, measurement being nothing else than a marking out and ascertaining of the thing measured: hence every art repudiates infinities. But individual existences are infinite, at least potentially.

\para[197] 6. The sixth argument is from the vileness of individual things. As the nobility of knowledge is weighed according to the nobility of the thing known, so the vileness also of the thing known seems to redound to the vileness of the knowledge. Therefore the excellent nobility of the divine mind does not permit of God knowing sundry most vile things that have individual existence.

\para[198] 7. The seventh argument is from the evil that is found in sundry individual things. Since the object known is in some manner in the knowing mind, and evil is impossible in God, it seems to follow that God can have no knowledge at all of evil and privation: only the mind that is in potentiality can know that, as privation can be only in potentiality.

\section*{Chapter 64. A list of things to be said concerning the Divine Knowledge}

\para[199] TO the exclusion of the above error we will show first that the divine mind does know individual things; secondly, that it knows things which actually are not; thirdly, that it knows future contingencies with infallible knowledge; fourthly, that it knows the motions of the will; fifthly, that it knows infinite things; sixthly, that it knows the vilest and least of things that be; seventhly, that it knows evils and all manner of privations or defects.

\section*{Chapter 65. That God Knows Individual Things}

\para[200] GOD knows things in so far as He is the cause of them. But the substantial effects of divine causation are individual things, universals not being substantial things, but having being only in individuals.

\para[201] 2. Since God's cognitive act is His essence, He must know all that is in any way in His essence; and as this essence is the first and universal principle of being and the prime origin of all, it virtually contains in itself all things that in any way whatsoever have being.

\para[202] 5. In the gradation of faculties it is commonly found that the higher faculty extends to more terms, and yet is one; while the range of the lower faculty extends to fewer terms, and even over them it is multiplied, as we see in the case of imagination and sense, for the single power of the imagination extends to all that the five senses take cognisance of, and to more. But the cognitive faculty in God is higher than in man: whatever therefore man knows by the various faculties of understanding, imagination and sense, God is cognisant of by His one simple intuition. God therefore is apt to know the individual things that we grasp by sense and imagination.

\para[203] 6. The divine mind, unlike ours, does not gather its knowledge from things, but rather by its knowledge is the cause of things; and thus its knowledge of things is a practical knowledge. But practical knowledge is not perfect unless it descends to individual cases: for the end of practical knowledge is work, which is done on individuals.

\para[204] 9. As the Philosopher argues against Empedocles, God would be very wanting in wisdom, if He did not know individual instances, which even men know.

\para[205] This truth is established also by the authority of Holy Scripture, for it is said: There is no creature invisible in his sight: also the contrary error is excluded by the text: Say not, I shall be hidden from God; and from the height of heaven who shall mind me? (Ecclus xvi, 16).

\para[206] From what has been said it is evident how the objection to the contrary (Chap. , i) is inconclusive: for though the mental presentation whereby divine understanding understands is immaterial, it is still a type both of matter and form, as being the prime productive principle of both.

\section*{Chapter 66. That God knows things which are not *}

\para[207] THE knowledge of the divine mind stands to things as the knowledge of the artificer to the products of his art. But the artificer by the knowledge of his art knows even those products of it which are not yet produced.

\para[208] 3. God knows other things besides Himself by His essence, inasmuch as His essence is the type of other things that come forth from Him (Chap. ). But since the essence of God is infinitely perfect (Chap. ), while of every other thing the being and perfection is limited, it is impossible for the whole sum of other things to equal the perfection of the divine essence. Therefore the representative power of that essence extends to many more things than the things that are. As then God knows entirely the power and perfection of His essence, His knowledge reaches not only to things that are, but also to things that are not.

\para[209] 6. The understanding of God has no succession, as neither has His being: it is all together, ever abiding, which is the essential notion of eternity, whereas the duration of time extends by succession of before and after. The proportion of eternity to the whole duration of time is as the proportion of an indivisible point to a continuous surface, --- not of that indivisible point which is a term of the surface, and is not in every part of its continuous extent: for to such a point an instant of time bears resemblance; but of that indivisible point which lies outside of the surface, and yet co-exists with every part or point of its continuous extent: for since time does not run beyond motion, eternity, which is altogether beyond motion, is no function of time. Again, since the being of the eternal never fails, eternity is present to every time or instant of time. Some sort of example of this may be seen in a circle: for a point taken on the circumference does not coincide with every other point; but the centre, lying away from the circumference, is directly opposite to every point on the circumference. Whatever therefore is in any portion of time, co-exists with the eternal, as present to it, although in respect to another portion of time it be past or future. But nothing can co-exist in presence with the eternal otherwise than with the whole of it, because it has no successive duration. Whatever therefore is done in the whole course of time, the divine mind beholds it as present throughout the whole of its eternity; and yet it cannot be said that what is done in a definite portion of time has always been an existing fact. The conclusion is that God has knowledge of things that in the course of time as yet are not.

\para[210] By these reasons it appears that God has knowledge of nonentities. But all nonentities do not stand in the same regard to His knowledge. Things that neither are, nor shall be, nor have been, are known by God as possible to His power: hence He does not know them as being anywise in themselves, but only as being within the compass of divine power. These sort of things are said by some to be known by God with the 'knowledge of simple understanding' (notitia simplicis intelligentiae). But as for those things that are present, past, or future to us, God knows them as they are within the compass of His power; and as they are within the compass of their own several created causes; and as they are in themselves; and the knowledge of such things is called the 'knowledge of vision' (notitia visionis). For of the things that are not yet with us, God sees not only the being that they have in their causes, but also the being that they have in themselves, inasmuch as His eternity is indivisibly present to all time. We must remember that God knows the being of everything through His own essence: for His essence is representable by many things that are not, nor ever shall be, nor ever have been. That same essence is the type of the power of every cause, in virtue of which power effects pre-exist in their causes. Again the being of everything, that it has in itself, is modelled upon the being of the divine essence. Thus then God knows nonentities inasmuch as in some way they have being, either in the power of God, or in their (creature) causes, or in themselves. To this the authority of Holy Scripture also gives testimony: All things are known to the Lord our God before their creation; as also, after they are fully made, he regardeth all (Ecclus XXIII, 29): and, Before I formed thee in the womb, I knew thee (Jer. I, 5).

\section*{Chapter 67. That God knows Individual Contingent Events *}

\para[211] HENCE we may gather some inkling of how God has had an infallible knowledge of all contingent events from eternity, and yet they cease not to be contingent. For contingency is not inconsistent with certain and assured knowledge except so far as the contingent event lies in the future, not as it is present. While the event is in the future, it may not be; and thus the view of him who reckons that it still be may be mistaken: but once it is present, for that time it cannot but be. Any view therefore formed upon a contingent event inasmuch as it is present may be a certitude. But the intuition of the divine mind rests from eternity upon each and every [one] of the events that happen in the course of time, viewing each as a thing present. There is nothing therefore to hinder God from having from eternity an infallible knowledge of contingent events.

\para[212] 2. A contingent event differs from a necessary event in point of the way in which each is contained in its cause. A contingent event is so contained in its cause as that it either may not or may ensue therefrom: whereas a necessary event cannot but ensue from its cause. But as each of these events is in itself, the two do not differ in point of reality; and upon reality truth is founded. In a contingent event, considered as it is in itself, there is no question of being or not being, but only of being: although, looking to the future, a contingent event possibly may not come off. But the divine mind knows things from eternity, not only in the being which they have in their causes, but also in the being which they have in themselves.

\para[213] 3. As from a necessary cause the effect follows with certainty, with like certainty does it follow from a contingent cause, when the cause is complete, provided no hindrance be placed. But as God knows all things (Chap. ). He knows not only the causes of contingent events, but like-wise the means whereby they may be hindered from coming off. He knows therefore with certitude whether they are going to come off or not.

\para[214] 6. The knowledge of God would not be true and perfect, if things did not happen in the way that God apprehends them to happen. But God, cognisant as He is of all being of which He is the principle, knows every event, not only in itself, but also in its dependence on any proximate causes on which it happens to depend: but the dependence of contingent events upon their proximate causes involves their ensuing upon them contingently. God therefore knows sundry events to happen, and to happen contingently: thus the certitude and truth of divine knowledge does not remove the contingency of events.

\para[215] 7. When it is said, 'God knows, or knew, this coming event,' an intervening medium is supposed between the divine knowledge and the thing known, to wit, the time to which the utterance points, in respect to which that which is said to be known by God is in the future. But really it is not in the future in respect of the divine knowledge, which existing in the instant of eternity is present to all things. In respect of such knowledge, if we set aside the time of speaking, it is impossible to say that so-and-so is known as non-existent; and the question never arises as to whether the thing possibly may never occur. As thus known, it should be said to be seen by God as already present in its existence. Under this aspect, the question of the possibility of the thing never coming to be can no longer be raised: what already is, in respect of that present instant cannot but be. The fallacy then arises from this, that the time at which we speak, when we say 'God knows,' co-exists with eternity; or again the last time that is marked when we say 'God knew'; and thus a relation of time, past or present, to future is attributed to eternity, which attribution does not hold; and thus we have fallacia accidentis.

\para[216] 8. Since everything is known by God as seen by Him in the present, the necessity of that being true which God knows is like the necessity of Socrates's sitting from the fact of his being seen seated. This is not necessary absolutely, 'by necessity of the consequent,' as the phrase is, but conditionally, or 'by necessity of the consequence.' For this conditional proposition is necessary: 'He is sitting, if he is seen seated.' Change the conditional proposition into a categorical of this form: 'What is seen sitting, is necessarily seated': it is clear that the proposition is true as a phrase, where its elements are taken together (compositam), but false as a fact, when its elements are separated (divisam). All these objections against the divine knowledge of contingent facts are fallacia compositionis et divisionis.

\para[217] That God knows future contingencies is shown also by the authority of Holy Scripture: for it is said of Divine Wisdom, It knows signs and portents beforehand, and the issues of times and ages (Wisd. viii, 8): and, There is nothing hidden from his eyes: from age to age he regardeth (Ecclus xxxix, 24, 25).

\section*{Chapter 68. That God knows the Motions of the Will}

\para[218] GOD knows the thoughts of minds and the volitions of hearts in virtue of their cause, as He is Himself the universal principle of being. All that in any way is, is known by God in His knowledge of His own essence (Chap. ). Now there is a certain reality in the soul, and again a certain reality in things outside the soul. The reality in the soul is that which is in the will or thought. God knows all these varieties of reality.

\para[219] 3. As God by knowing His own being knows the being of everything, so by knowing His own act of understanding and will He knows every thought and volition.

\para[220] 5. God knows intelligent substances not less well than He knows or we know sensible substances, seeing that intelligent substances are more knowable, as being better actualised. This is confirmed by the testimony of Holy Scripture: --- God searcher of hearts and reins (Ps. vii, 10): Hell and perdition are before the Lord: how much more the hearts of the sons of men? (Prov. xi, 11): He needed not that any one should bear testimony of what was in man: for he himself knew what was in man (John ii, 25).

\para[221] The dominion of the will over its own acts, whereby it has it in its power to will and not to will, is inconsistent with will-force being determined to one fixed mode of action: it is inconsistent also with the violent interference of any external agency; but it is not inconsistent with the influence of that Higher Cause, from whence it is given to the will both to be and to act. And thus in the First Cause, that is, in God, there remains a causal influence over the motions of the will, such that, in knowing Himself, God is able to know these motions.

\section*{Chapter 69. That God knows infinite things *}

\para[222] BY knowing Himself as the cause of things, He knows things other than Himself (Chap. ). But He is the cause of infinite things, if beings are infinite, for He is the cause of all things that are.

\para[223] 2. God knows His own power perfectly (Chap. , ). But power cannot be perfectly known, unless all the objects to which it extends are known, since according to that extent the amount of the power may be said to be determined. But His power being infinite (Chap. ) extends to things infinite, and therefore also His knowledge.

\para[224] 3. If the knowledge of God extends to all things that in any sort of way are, He must not only know actual being, but also potential being. But in the physical world there is potential infinity, though not actual infinity, as the Philosopher proves. God therefore knows infinite things, in the way that unity, which is the principle of number, would know infinite species of number if it knew whatever is in its potentiality: for unity is in promise and potency every number.

\para[225] 4. God in His essence, as in a sort of exemplar medium, knows other things. But as He is a being of infinite perfection, there can be modelled upon Him infinite copies with finite perfections, because no one of these copies, nor any number of them put together, can come up to the perfection of their exemplar; and thus there always remains some new way for any copy taken to imitate Him.

\para[226] 10. The infinite defies knowledge in so far as it defies counting. To count the parts of the infinite is an intrinsic impossibility, as involving a contradiction. To know a thing by enumeration of its parts is characteristic of a mind that knows part after part successively, not of a mind that comprehends the several parts together. Since then the divine mind knows all things together without succession, it has no more difficulty in knowing things infinite than in knowing things finite.

\para[227] 11. All quantity consists in a certain multiplication of parts; and therefore number is the first of quantities. Where then plurality makes no difference, no difference can be made there by anything that follows upon quantity. But in God's knowledge many things are known in one, not by many different presentations, but by that one species, or presentation, which is the essence of God. Hence a multitude of things is known by God all at once; and thus plurality makes no difference in God's knowledge: neither then does infinity, which follows upon quantity.

\para[228] In accordance with this is what is said in Psalm cxlvi: And of his wisdom there is no telling.

\para[229] From what has been said it is clear why our mind does not know the infinite as the divine mind does. Our mind differs from the divine mind in four respects; and they make all the difference. The first is that our mind is simply finite, the divine mind infinite. The second is that as our mind knows different things by different impressions, it cannot extend to an infinity of things, as the divine mind can. The third results in this way, that as our mind is cognisant of different things by different impressions, it cannot be actually cognisant of a multitude of things at the same time; and thus it could not know an infinity of things except by counting them in succession, which is not the case with the divine mind, which discerns many things at once as seen by one presentation. The fourth thing is that the divine mind is cognisant of things that are and of things that are not (Chap. ).

\para[230] It is also clear how the saying of the Philosopher, that the infinite, as infinite, is unknowable, is in no opposition with the opinion now put forth: because the notion of infinity attaches to quantity; consequently, for infinite to be known as infinite, it would have to be known by the measurement of its parts, for that is the proper way of knowing quantity: but God does not know the infinite in that way. Hence, so to say, God does not know the infinite inasmuch as it is infinite, but inasmuch as, to His knowledge, it is as though it were finite.

\para[231] It is to be observed however that God does not know an infinity of things with the 'knowledge of vision,' because infinite things neither actually are, nor have been, nor shall be, since, according to the Catholic faith, there are not infinite generations either in point of time past or in point of time to come. But He does know an infinity of things with the 'knowledge of simple understanding': for He knows infinite things that neither are, nor have been, nor shall be, and yet are in the power of the creature; and He also knows infinite things that are in His own power, which neither are, nor shall be, nor have been. Hence to the question of the knowledge of particular things it may be replied by denial of the major: for particular things are not infinite: if however they were, God would none the less know them.

\section*{Chapter 70. That God knows Base and Mean Things *}

\para[232] THE stronger an active power is, to the more remote objects does it extend its action. But the power of the divine mind in knowing things is likened to active power: since the divine mind knows, not by receiving aught from things, but rather by pouring its influence upon things. Since then God's mind is of infinite power in understanding (Chap. ), its knowledge must extend to the remotest objects. But the degree of nobility or baseness in all things is determined by nearness to or distance from God, who is the fulness of nobility. Therefore the very vilest things in being are known to God on account of the exceeding great power of His understanding.

\para[233] 2. Everything that is, in so far as it has place in the category of substance or quality, is in actuality: it is some sort of likeness of the prime actuality, and is ennobled thereby. Even potential being, from its reference to actuality shares in nobility, and so comes to have the name of 'being.' It follows that every being, considered in itself, is noble; and is only mean and vile in comparison with some other being, nobler still. But the noblest creatures are removed from God at a distance not less than that which separates the highest in the scale of creation from the lowest. If then the one distance were to bar God's knowledge, much more would the other; and the consequence would be that God would know nothing beyond Himself.

\para[234] 3. The good of the order of the universe is nobler than any part of the universe. If then God knows any other noble nature, most of all must He know the order of the universe. But this cannot be known without taking cognisance at once of things nobler and things baser: for in the mutual distances and relations of these things the order of the universe consists.

\para[235] 4. The vileness of the objects of knowledge does not of itself redound on to the knower; for it is of the essence of knowledge that the knower should contain within himself impressions of the object known according to his own mode and manner. Accidentally however the vileness of the objects known may redound upon the knower, either because in knowing base and mean things he is withdrawn from the thought of nobler things, or because from the consideration of such vile objects he is inclined to some undue affections: which cannot be the case with God.

\para[236] 5. A power is not judged to be small, which extends to small things, but only that which is limited to small things. A knowledge therefore that ranges alike over things noble and things mean, is not to be judged mean; but that knowledge is mean, which ranges only over mean things, as is the case with us: for we make different studies of divine and of human things, and there is a different science of each. But with God it is not so; for with the same knowledge and the same glance He views Himself and all other beings.

\para[237] With this agrees what is said of the Divine Wisdom: It findeth place everywhere on account of its purity, and nothing defiled stealeth in to corrupt it (Wisdom vii, 24, 25).

\section*{Chapter 71. That God knows Evil Things}

\para[238] WHEN good is known, the opposite evil is known. But God knows all particular good things, to which evil things are opposed: therefore God knows evil things.

\para[239] 2. The ideas of contraries, as ideas in the mind, are not contrary to one another: otherwise they could not be together in the mind, or be known together: the idea therefore whereby evil is known is not inconsistent with good, but rather belongs to the idea of good (ratio qua cognoscitur malum ad rationem boni pertinet). If then in God, on account of His absolute perfection, there are found all ideas of goodness (rationes bonitatis, as has been proved (Chap. ), It follows that there is in Him the idea (ratio) whereby evil is known.

\para[240] 3. Truth is the good of the understanding: for an understanding is called good inasmuch as it knows the truth. But truth is not only to the effect that good is good, but also that evil is evil: for as it is true that what is, is, so it is true that what is not, is not. The good of the understanding therefore consists even in the knowledge of evil. But since the divine understanding is perfect in goodness, there cannot be wanting to it any of the perfections of understanding; and therefore there is present to it the knowledge of things evil.

\para[241] 4. God knows the distinction of things (Chap. ). But in the notion of distinction there is negation: for those things are distinct, of which one is not another: hence the first things that are of themselves distinct, mutually involve the exclusion of one another, by reason of which fast negative propositions are immediately verified of them, e.g., 'No quantity is a substance.' God then knows negation. But privation is a sort of negation: He therefore knows privation, and consequently evil, which is nothing else than a privation of due perfection.

\para[242] 8. In us the knowledge of evil things is never blameworthy in mere point of knowledge, that is in the judgement that is passed about evil things, but accidentally, inasmuch as by the observation of evil things one is sometimes inclined to evil. But that cannot be in God; and therefore there is nothing to prevent His knowing evil.

\para[243] With this agrees what is said, that Evil surpasseth not [God's] wisdom (Wisd. vii 30) and, Hell and perdition are before the Lord (Prov. xv, 11) and, My offences are not hidden from thee (Ps. lxviii, 6); and, He knoweth the vanity of men, and seeing doth he not consider iniquity? (Job xi, 11.)

\para[244] It is to be observed however that if God's knowledge were so limited as that His knowledge of Himself did not involve His knowing other beings of finite and partial goodness, at that rate He would nowise know privation or evil: because to the good which is God Himself there is no privation opposed, since privation and its opposite are naturally about the same object; and so to that which is pure actuality no privation is opposed, and consequently no evil either. Hence on the supposition that God knows Himself alone, by knowing the excellences of His own being, He will not know evil. But because in knowing Himself He knows beings in which privations naturally occur, He must know the opposite privations, and the evils opposite to particular goods.

\para[245] It must be further observed that as God, without any argumentative process, knows other beings by knowing Himself, so there is no need of His knowledge being argumentative in coming to the knowledge of evil things through good things: for good is as it were the ground of the knowledge of evil, evil being nothing else than privation of good: hence what is evil is known through what is good as things are known through their definitions, not as conclusions through their premises.

\section*{Chapter 72. That God has a Will}

\para[246] FROM the fact that God has understanding, it follows that He has a will. Since good apprehended in understanding is the proper object of the will, understood good, as such, must be willed good. But anything understood involves an understanding mind. A mind then that understands good, must, as such, be a mind that wills good.

\para[247] 3. What is consequent upon all being, is a property of being, as such. Such a property must be found in its perfection in the first and greatest of beings. Now it is a property of all being to seek its own perfection and the preservation of its own existence. Every being does this in its own way: intelligent beings, by their will: animals, by their sensitive appetite: unconscious nature, by a certain physical nisus. It makes a difference however whether the thing craved for is possessed or not. Where it is not possessed, the nisus of desire proper to each several kind goes out to seek what is wanting: where the thing is possessed, it is rested in and clung to. This characteristic of all being cannot be wanting in the first of beings, which is God. Since then God has understanding, He has also a will, whereby He takes complacency in His own being and His own goodness.

\para[248] 4. The more perfect the act of understanding is, the more delightful to the understanding mind. But God has understanding and a most perfect act thereof (Chap. ): therefore that act yields Him the utmost delight. But as sensible delight is through the concupiscible appetite, so is intellectual delight through the will. God then has a will. This will of God the testimonies of Holy Scripture confess: All things whatsoever he hath willed, the Lord hath done (Ps. cxxxiv, 6): Who resisteth his will? (Rom. ix, 19).

\section*{Chapter 73. That the Will of God is His Essence}

\para[249] GOD has will inasmuch as He has understanding. But He has under- standing by His essence (Chap. , ), and therefore will in like manner.

\para[250] 2. The act of will is the perfection of the agent willing. But the divine being is of itself most perfect, and admits of no superadded perfection (Chap. ): therefore in God the act of His willing is the act of His being.

\para[251] 3. As every agent acts inasmuch as it is in actuality, God, being pure actuality, must act by His essence. But to will is an act of God: therefore God must will by His essence.

\para[252] 4. If will were anything superadded to the divine substance, that substance being complete in being, it would follow that will was something adventitious to it as an accident to a subject; also that the divine substance stood to the divine will as potentiality to actuality; and that there was composition in God: all of which positions have been rejected (Chap. , , ).

\section*{Chapter 74. That the Object of the Will of God in the First Place is God Himself}

\para[253] GOOD understood is the object of the will. But what is understood by God in the first place is the divine essence: therefore the divine essence is the first object of the divine will.

\para[254] 3. The object in the first place willed is the cause of willing to every willing agent. For when we say, 'I wish to walk for the benefit of my health,' we consider that we are assigning a cause; and if we are further asked, 'Why do you wish to benefit your health?' we shall go on assigning causes until we come to the final end, which is the object willed in the first place, and is in itself the cause of all our willing. If then God wills anything else than Himself in the first place, it will follow that that 'something else' is to Him a cause of willing. But His willing is His being (Chap. ), Therefore something else will be the cause of His being, which is contrary to the notion of the First Being.

\section*{Chapter 75. That God in willing Himself wills also other things besides Himself *}

\para[255] EVERY one desires the perfection of that which for its own sake he wills and loves: for the things which we love for their own sakes we wish to be excellent, and ever better and better, and to be multiplied as much as possible. But God wills and loves His essence for its own sake. Now that essence is not augmentable and multipliable in itself (Chap. ), but can be multiplied only in its likeness, which is shared by many. God therefore wills the multitude of things, inasmuch as He wills and loves His own perfection.

\para[256] 3. Whoever loves anything in itself and for itself, wills consequently all things in which that thing is found: as he who loves sweetness in itself must love all sweet things. But God wills and loves His own being in itself and for itself; and all other being is a sort of participation by likeness of His being.

\para[257] 6. The will follows the understanding. But God with His understanding understands Himself in the first place, and in Himself understands all other things: therefore in like manner He wills Himself in the first place, and in willing Himself wills all other things.

\para[258] This is confirmed by the authority of Holy Scripture: Thou lovest all things that are, and hatest nothing of the things that thou hast made (Wisd. xi, 2)

\section*{Chapter 76. That with one and the same Act of the Will God wills Himself and all other Beings}

\para[259] EVERY power tends by one and the same activity to its object and to that which makes the said object an object to such a power, as with the same vision we see light and the colour which is made actually visible by light. But when we wish a thing for an end, and for that alone, that which is desired for the end receives from the end its character of an object of volition. Since then God wills all things for Himself (Chap. ), with one act of will He wills Himself and other things.

\para[260] 2. What is perfectly known and desired, is known and desired to the whole extent of its motive power. But a final end is a motive not only inasmuch as it is desired in itself, but also inasmuch as other things are rendered desirable for its sake. He therefore who perfectly desires an end, desires it in both these ways. But it is impossible to suppose any volitional act of God, by which He should will Himself, and not will Himself perfectly: since there is nothing imperfect in God. By every act therefore by which He wills Himself, He wills Himself and other things for His own sake absolutely; and other things besides Himself He does not will except inasmuch as He wills Himself.

\para[261] 3. As promises are to conclusions in things speculative, so is the end to the means in things practical and desirable: for as we know conclusions by premises, so from the end in view proceeds both the desire and the carrying out of the means. If then one were to wish the end apart, and the means apart, by two separate acts, there would be a process from step to step in his volition (Chap. ). But this is impossible in God, who is beyond all movement.

\para[262] 7. To will belongs to God inasmuch as He has understanding. As then by one act He understands Himself and other beings, inasmuch as His essence is the pattern of them all, so by one act He wills Himself and all other beings, inasmuch as His goodness is the type of all goodness.

\section*{Chapter 77. That the Multitude of the Objects of God's Will is not inconsistent with the Simplicity of His Substance}

\para[263] GOD wills other things inasmuch as He wills His own goodness (Chap. ). Things then come under the will of God according as they are included in His goodness. But in His goodness all things are one: for they are in Him according to the mode that befits Him; material things, immaterially; and things many, in union (Chap. , ). Thus the multitude of the objects of the divine will does not multiply the divine substance.

\section*{Chapter 78. That the Divine Will reaches to the good of Individual Existences}

\para[264] THE excellence of order in the universe appears in two ways, first, inasmuch as the whole universe is referred to something beyond the universe, as an army to its leader: secondly, inasmuch as the parts of the universe are referred to one another, like the parts of an army; and the second order is for the sake of the first. But God, in willing Himself as an end, wills other things in their reference to Him as an end. He wills therefore the excellence of order in the universe in reference to Himself, and the excellence of order in the universe in mutual reference of its parts to one another. But the excellence of order is made up of the good of individual existences.

\para[265] This is confirmed by the authority of Scripture: God saw the light, that it was good (Gen. i, 4); and similarly of His other works; and lastly of them altogether: God saw all things that he had made, and they were very good (Gen. i, 31).

\section*{Chapter 79. That God wills things even that as yet are not}

\para[266] SOME one might perhaps think that God wills only the things that are: for correlatives go together; and if one perishes, the other perishes; if then willing supposes a relation of the willing subject to the object willed, none can will any but things that are. Besides, the will and its objects are to one another as Creator and creature: now God cannot be called Creator, or Lord, or Father, except of things that are: neither then can He be said to will any but things that are. And it may be further argued, that if the divine will is invariable, as is the divine being, and wills only actual existences, it wills nothing but what always is.

\para[267] Let us say then in answer to these objections, that as good apprehended by the intellect moves the will, the act of the will must follow the condition of the mental apprehension. Now the mind apprehends the thing, not only as it is in the mind, but also as it is in its own nature: for we not only know that the thing is understood by us (for that is the meaning of its being 'in the mind'), but also that the thing exists, or has existed, or is to exist in its own nature. Though then at the time the thing has no being other than in the mind, still the mind stands related to it, not as it is in the mind, but as it is in its own nature, which the mind apprehends. Therefore the relation of the divine will to a non-existent thing is to the thing according as it is in its own nature, attached to some certain time, and not merely to the thing as it is in the knowledge of God. For God wills the thing, that is not now, to be in some certain time: He does not merely will it inasmuch as He Himself understands it. Nor is the relation of the will to its object similar to the relation of Creator to creature, of Maker to made, of Lord to subject. For will, being an immanent act, does not involve the actual external existence of the thing willed: whereas making and creating and governing do signify an action terminated to an external effect, such that without its existence such action is unintelligible.

\section*{Chapter 80. That God of necessity wills His own Being and His own Goodness}

\para[268] GOD wills His own being and His own goodness as His first object and reason for willing all other things (Chap. ), and this He wills in everything that He does will. Nor is it possible for Him to will it merely potentially: He must will it actually, as His willing is His being.

\para[269] 4. All things, in so far as they have existence, are likened to God, who is the first and greatest being. But all things, in so far as they have existence, cherish their own being naturally in such manner as they can. Much more therefore does God cherish His own being naturally.

\section*{Chapter 81. That God does not of necessity love other things than Himself}

\para[270] A WILL does not of necessity tend to the means to an end, if the end can be had without those means. Since then the divine Goodness can be without other beings, --- nay, other beings make no addition to it, --- God is under no necessity of willing other things from the fact of His willing His own goodness.

\para[271] 2. Since good, understood to be such, is the proper object of the will, the will may fasten on any object conceived by the intellect in which the notion of good is fulfilled. Hence though the being of anything, as such, is good, and its not-being, as such, is evil; still the very not- being of a thing may become an object to the will, though not of necessity, by reason of some notion of good fulfilled: for it is good for a thing to be, even though some other thing is not. The only good then which the will by the terms of its constitution cannot wish not to be, is the good whose non-existence would destroy the notion of good altogether. Such a good is no other than God. The will then by its constitution can will the non-existence of anything else except of God. But in God there is will according to the fulness of the power of willing. God then can will the non-existence of any other being besides Himself.

\para[272] 3. God in willing His own goodness wills also other things than Himself as sharing His goodness. But since the divine goodness is infinite, and partakable in infinite ways, if by the willing of His own goodness He of necessity willed the beings that partake of it, the absurdity would follow that He must will the existence of infinite creatures sharing His goodness in infinite ways: because, if He willed them, those creatures would exist, since His will is the principle of being to creatures. We must consider therefore why God of necessity knows other beings than Himself, and yet does not of necessity will them to exist, notwithstanding that His understanding and willing of Himself involves His understanding and willing other beings. The reason of it is this: an intelligent agent's understanding anything arises from a certain condition of the understanding, --- for by a thing being actually understood its likeness is in the mind: but a volitional agent's willing anything arises from a certain condition of the object willed, --- for we will a thing either because it is an end, or because it is a means to an end. Now the divine perfection necessarily requires that all things should so be in God as to be understood in Him. But the divine goodness does not of necessity require that other things should exist to be referred to Him as means to an end; and therefore it is necessary that God should know other things, but not that He should will other things. Hence neither does He will all things that are referable to His goodness: but He knows all things which are in any way referable to His essence, whereby He understands.

\section*{Chapter 82. Arguments against the aforesaid Doctrine and Solutions of the same}

\para[273] THESE awkward consequences seem to follow, if any things that God wills He does not will of necessity.

\para[274] 1. If the will of God in respect of certain objects of will is not determined by any of them, it seems to be indifferent. But every faculty that indifferent is in a manner in potentiality.

\para[275] 2. Since potential being, as such, is naturally changeable, --- for what can be can also not be, --- it follows that the divine will is variable.

\para[276] 4. Since what hangs loose, indifferent between two alternatives, does not tend to one rather than to the other, unless it be determined by one or other, either God wills none of the things to which He is indifferent, or He is determined by one or other of them, in which case there must be something antecedent to God to determine Him.

\para[277] But none of the above objections can stand.

\para[278] 1. The indifference, or indeterminateness, of a faculty may be attributable either to the faculty itself or to its object. To the faculty itself, when its indeterminateness comes from its not having yet attained to its perfection. This argues imperfection in the faculty, and an unfulfilled potentiality, as we see in the mind of a doubter, who has not yet attained to premises sufficient to determine him to take either of two sides. To the object of the faculty, when the perfect working of the faculty does not depend on its adoption of either alternative, and yet either alternative may be adopted, as when art may employ different instruments to do the same work equally well. This argues no imperfection in the faculty, but rather its pre-eminent excellence, inasmuch as it rises superior to both opposing alternatives, and therefore is indifferent to both and determined by neither. Such is the position of the divine will with respect to things other than itself. Its perfection depends on none of them; being as it is intimately conjoined with its own last end and final perfection.

\para[279] 2. In the divine will there is no potentiality. Unnecessitated, it prefers one alternative to another respecting the creatures which it causes to be. It is not to be looked upon as being in a potential attitude to both alternatives, so as first to be potentially willing both, and then to be actually willing one. It is for ever actually willing whatever it wills, as well its own self as the creatures which are the objects of its causation. But whatever creature God wills to exist, that creature stands in no necessary relation to the divine goodness, which is the proper object of the divine will.

\para[280] 4. We cannot admit that either the divine will wills none of the effects of its causation, or that its volition is determined by some exterior object. The proper object of the will is good apprehended as such by the understanding. Now the divine understanding apprehends, not only the divine being, or divine goodness, but other good things likewise (Chap. ); and it apprehends them as likenesses of the divine goodness and essence, not as constituent elements of the same. Thus the divine will tends to them as things becoming its goodness, not as things necessary to its goodness. So it happens also in our will: which, when it inclines to a thing as absolutely necessary to its end, tends to it with a certain necessity; but when it tends to a thing solely on account of its comeliness and appropriateness, does not tend to it necessarily.

\section*{Chapter 83. That God wills anything else than Himself with an Hypothetical Necessity *}

\para[281] IN every unchangeable being, whatever once is, cannot afterwards cease to be. Since then God's will is unchangeable, supposing Him to will anything, He cannot on that supposition not will it.

\para[282] 2. Everything eternal is necessary. But God's will for the causation of any effect is eternal: for, as His being, so His willing is measured by eternity. That will therefore is necessary, yet not absolutely so, since the will of God has no necessary connexion with this objection willed. It is therefore necessary hypothetically, on a supposition.

\para[283] 3. Whatever God once could do, He can still. His power does not grow less, as neither does His essence. But He cannot now not-will what He is already supposed to have willed, because His will cannot change: therefore He never could not-will whatever He once willed (nunquam potuit non velle quidquid voluit). It is therefore hypothetically necessary for Him to have willed whatever He has willed, as it is for Him to will whatever He does will: but in neither case is the necessity absolute.

\para[284] 4. Whoever wills anything, necessarily wills all that is necessarily requisite to that purpose, unless there be some defect on his part, either by ignorance, or because his will sometimes is drawn away by some passion from a right choice of means to the end: nothing of which can be said of God. If God then in willing Himself wills anything else besides Himself, He needs must will all that is necessarily required to the effecting of the thing willed, as it is necessary that God should will the being of a rational soul, if He wills the being of a man.

\section*{Chapter 84. That the Will of God is not of things in themselves Impossible}

\para[285] THOSE things are in themselves impossible, which involve an inconsistency, as that man should be an ass, which involves the rational being irrational. But what is inconsistent with a thing, excludes some one of the conditions requisite to it, as being an ass excludes a man's reason. If therefore God necessarily wills the things requisite to that which by supposition He does will, it is impossible for Him to will what is inconsistent therewith.

\para[286] 2. God, in willing His own being, wills all other things, that He does will, in so far as they have some likeness to it. But in so far as anything is inconsistent with the notion of being as such, there cannot stand therein any likeness to the first or divine being, which is the fountain of being. God therefore cannot will anything that is inconsistent with the notion of being as such, as that anything should be at once being and not being, that affirmation and negation should be true together, or any other such essential impossibility, inconsistency, and implied contradiction.

\para[287] 3. What is no object of the intellect, can be no object of the will. But essential impossibilities, involving notions mutually inconsistent, are no objects of intellect, except perchance through the error of a mind that does not understand the proprieties of things, which cannot be said of God.

\section*{Chapter 85. That the Divine Will does not take away Contingency from things}

\para[288] HYPOTHETICAL necessity in the cause cannot lead to absolute necessity in the effect. But God's will about a creature is not absolutely necessary, but hypothetically so (Chap. ). Therefore the divine will is no argument of absolute necessity in creatures. But only this absolute necessity excludes contingency: for even a contingent fact may be extended either way into an hypothetical necessity: thus it is necessary that Socrates moves, if he runs. It does not therefore follow that a thing happens of necessity, if God wills it: all that holds is the necessary truth of this conditional: 'If God wills anything, the thing will be': but the 'consequent' (as distinguished from the 'consequence') need not be a necessary truth.

\section*{Chapter 86. That Reason can be assigned for the Divine Will *}

\para[289] THE end is a reason for willing the means. But God wills His own goodness as an end, and all things else as means thereto: His goodness therefore is a reason why He wills other things different from Himself.

\para[290] 2. The good of a part is ordained to the end of the good of the whole, as the imperfect to the perfect. But things become objects of the divine will according as they stand in the order of goodness. It follows that the good of the universe is the reason why God wills every good of any part of the universe.

\para[291] 3. Supposing that God wills anything, it follows of necessity that He wills the means requisite thereto. But what lays on others a necessity for doing a thing, is a reason for doing it. Therefore the accomplishment of a purpose, to which such and such means are requisite, is a reason to God for willing those means.

\para[292] We may therefore proceed as follows. God wishes man to have reason, to the end that he may be man: He wishes man to be, to the end of the completion of the universe: He wishes the good of the universe to be, because it befits His own goodness. The same proportion however is not observable in all three stages of this ratiocination. The divine goodness does not depend on the perfection of the universe, and receives no accession thereby. The perfection of the universe, though depending necessarily on the good of some particular components, which are essential parts of the universe, has no necessary dependence on others, although even from them some goodness or beauty accrues to the universe, such things serving solely for the fortification (munimentum) or embellishment of the rest. But any particular good depends absolutely on the elements that are requisite to it: and still even such goods have adjuncts that go merely to better their condition. Sometimes therefore the reason of the divine will involves mere becomingness, sometimes utility, sometimes also hypothetical necessity, but never absolute necessity, except when the object of God's volition is God Himself.

\section*{Chapter 87. That nothing can be a Cause to the Divine Will}

\para[293] THOUGH some reason may be assigned for the divine will, yet it does not follow that there is any cause of that will's volition. For the cause of volition is the end in view: now the end in view of the divine will is its own goodness: that then is God's cause of willing, which is also His own act of willing. But of other objects willed by God none is to God a cause of willing, but one of them is cause to another of its being referred to the divine goodness, and thus God is understood to will one for the sake of another. But clearly we must suppose no passing from point to point of God's will, where there is only one act, as shown above of the divine intellect (Chap. ). For God by one act wills His own goodness and all other things, as His action is His essence. By this and the previous chapter the error is excluded of some who say that all things proceed from God by sheer will, so that no reason is to be rendered of anything that He does beyond the fact that God so wills. Which position is even contrary to divine Scripture, which tells us that God has done all things according to the order of His wisdom: Thou hast done all things in wisdom (Ps. ciii, 24) ; and God has shed wisdom over all his works (Ecclus i,10).

\section*{Chapter 88. That there is a Free Will in God}

\para[294] GOD does not necessarily will things outside Himself (Chap. ).

\para[295] 3. Will is of the end: choice of the means. Since then God wills Himself as end, and other things as means, it follows that in respect of Himself He has will only, but in respect of other things choice. But choice is always an act of free will.

\para[296] 4. Man by free will is said to be master of his own acts. But this mastery belongs most of all to the Prime Agent, whose act depends on no other.

\section*{Chapter 89. That there are no Passions in God}

\para[297] PASSION is not in the intellectual appetite, but only in the sensitive. But in God there is no sensitive appetite, as there is no sensible knowedge.

\para[298] 2. Every passion involves some bodily alteration, a thing impossible in the incorporeal Deity.

\para[299] 3. In every passion the subject is more or less drawn out of his essential condition or connatural disposition: which is not possible in the unchangeable God.

\para[300] 4. Every passion fixes determinedly on some one object, according to the mode and measure of the passion. Passion, like physical nature, rushes blindly at some one thing: that is why passion needs repressing and regulating by reason. But the divine will is not determined of itself to any one object in creation: but proceeds according to the order of its wisdom(Chap. ).

\para[301] 5. Every passion is the passion of a subject that is in potentiality. But God is altogether free from potentiality, being pure actuality.

\para[302] Thus every passion, generically as such, is removed from God. But certain passions are removed from God, not only generically, but also specifically. For every passion takes its species from its object: if then an object is altogether unbefitting for God, the passion specified by that object is removed from God also on specific grounds. Such a passion is Sadness and Grief, the object of which is evil already attaching to the sufferer. Hope, again, though it has good for its object, is not of good obtained, but to be obtained, a relation to good which is unbefitting for God by reason of His so great perfection, to which addition is impossible. Much more does that perfection exclude any potentiality in the way of evil. But Fear regards an evil that may be imminent. In two ways then Fear, specifically as such, is removed from God, both because it supposes a subject that is in potentiality, and because it has for its object some evil that may come to be in the subject. Regret again, or Repentance, is repugnant to God, as well because it is a species of sadness, as also because it involves a change of will.

\para[303] Moreover, without an error of the intellectual faculty, it is impossible for good to be mistaken for evil. And only in respect of private advantages is it possible for the loss of one being to be the gain of another. But to the general good nothing is lost by the good of any private member; but every private good goes to fill in the public good. But God is the universal good, by partaking in whose likeness all other things are called good. No other being's evil then can possibly be good for God. Nor again, seeing that God's knowledge makes no mistakes, can He apprehend as evil that which is simply good, and no evil to Him. Envy therefore is impossible to God, specifically as Envy, not only because it is a species of sadness, but also because it is sadness at the good of another, and thus takes the good of another as evil to itself.

\para[304] It is part of the same procedure to be sad at good and to desire evil. Such sadness arises from good being accounted evil: such desire, from evil being accounted good. Now Anger is desire of the evil of another for vengeance' sake. Anger then is far from God by reason of its species, not only because it is a species of sadness, but also because it is a desire of vengeance, conceived for sadness at an injury done one.

\section*{Chapter 90. That there is in God Delight and Joy}

\para[305] THERE are some passions which, though they do not befit God as passions, nevertheless, so far as their specific nature is considered, do not involve anything inconsistent with divine perfection. Of the number of these is Delight and Joy. Joy is of present good. Neither by reason of its object, which is good, nor by reason of the relation in which the object, good actually possessed, stands to the subject, does joy specifically contain anything inconsistent with divine perfection. Hence it is manifest that joy or Delight has being properly in God. For as good and evil apprehended is the object of the sensitive appetite, so also is it of the intellectual appetite, or will. It is the ordinary function of both appetites to pursue good and to shun evil, either real or apparent, except that the object of the intellectual appetite is wider than that of the sensitive, inasmuch as the intellectual appetite regards good and evil simply, while the sensitive appetite regards good and evil felt by sense; as also the object of intellect is wider than the object of sense. But the activities of appetite are specified by their objects. There exist therefore in the intellectual appetite, or will, activities specifically similar to the activities of the sensitive appetite, and differing only in this, that in the sensitive appetite they are passions on account of the implication of a bodily organ, but in the intellectual appetite they are simple activities. For as by the passion of fear, coming over the sensitive appetite, one shuns evil looming in the future, so the intellectual appetite works to the same effect without passion. Since then joy and Delight are not repugnant to God specifically, but only inasmuch as they are passions, it follows that they are not wanting even in the divine will.

\para[306] 2. Joy and Delight are a sort of rest of the will in its object. But God singularly rests in Himself as in the first object of His own will, inasmuch as He has all sufficiency in Himself.

\para[307] 3. Delight is the perfection of activity, perfecting activity as bloom does youth. But the activity of the divine understanding is most perfect. If therefore our act of understanding, coming to its perfection, yields delight, most delightful must be the act whereby God understands.

\para[308] 4. Everything naturally feels joy over what is like itself, except accidentally, inasmuch as the likeness hinders one's own gain, and 'two of a trade' quarrel. But every good thing is some likeness of the divine goodness, and nothing is lost to God by the good of His creature. Therefore God rejoices in good everywhere.

\para[309] Joy and Delight differ in our consideration: for Delight arises out of good really conjoined with the subject; while Joy does not require this real conjunction, but the mere resting of the will on an agreeable object is sufficient for it. Hence, strictly speaking, Delight is at good conjoined with the subject: Joy over good external to the subject. Thus, in strict parlance, God takes delight in Himself: but has Joy both over Himself and over other things.

\section*{Chapter 91. That there is Love in God. *}

\para[310] IT is of the essential idea of love, that whoever loves wishes the good of the object loved. But God wishes His own good and the good of other beings (Chap. ); and in this respect He loves Himself and other beings.

\para[311] 2. It is a requisite of true love to love the good of another inasmuch as it is his good. But God loves the good of every being as it is the good of that being, though He does also subordinate one being to the profit of another.

\para[312] 3. The essential idea of love seems to be this, that the affection of one tends to another as to a being who is in some way one with himself. The greater the bond of union, the more intense is the love. And again the more intimately bound up with the lover the bond of union is, the stronger the love. But that bond whereby all things are united with God, namely, His goodness, of which all things are imitations, is to God the greatest and most intimate of bonds, seeing that He is Himself His own goodness. There is therefore in God a love, not only true, but most perfect and strong.

\para[313] But some might be of opinion that God does not love one object more than another; for a higher and a lower degree of intensity of affection is characteristic of a changeable nature, and cannot be attributed to God, from whom all change is utterly removed. Besides, wherever else there is mention of any divine activity, there is no question of more and less: thus one thing is not known by God more than another. In answer to this difficulty we must observe that whereas other activities of the soul are concerned with one object only, love alone seems to tend to two. For love wishes something to somebody: hence the things that we desire, we are properly said to 'desire,' not to 'love,' but in them we rather love ourselves for whom we desire them. Every divine act then is of one and the same intensity; but love may be said to admit of 'greater and less' in two ways, either in point of the good that we will to another, in which way we are said to love him more to whom we wish greater good; or again in point of the intensity of the act, in which way we are said to love him more to whom we wish, not indeed a greater good, but an equal good more fervently and effectually. In the former way then there is nothing to object to in the saying that God loves one more than another, inasmuch as He wishes him a greater good: but, understood of the second way, the saying is not tenable.

\para[314] Hence it appears that of our affections there is none that can properly be in God except joy and love, though even these are in Him not by way of passion, as they are in us. That there is in God joy or delight is confirmed by the authority of Holy Scripture. I was delighted day by day playing before him, says the Divine Wisdom, which is God (Prov. viii, 30). The Philosopher also says that God ever rejoices with one simple delight. The Scripture also speaks of love in God: With everlasting love I have loved thee (Jer. xxxi, 3); For the Father himself loveth you (John xvi, 27).

\para[315] But even other affections (affectiones), which are specifically inconsistent with divine perfection, are predicated in Holy Writ of God, not properly but metaphorically, on account of likeness of effects. Thus sometimes the will in following out the order of wisdom tends to the same effect to which one might be inclined by a passion, which would argue a certain imperfection: for the judge punishes from a sense of justice, as an angry man under the promptings of anger. So sometimes God is said to be 'angry,' inasmuch as in the order of His wisdom He means to punish some one: When his anger shall blaze out suddenly (Ps. ii, 13). He is said to be 'compassionate,' inasmuch as in His benevolence He takes away the miseries of men, as we do the same from a sentiment of pity: The Lord is merciful and compassionate, patient and abounding in mercy (Ps. cli, 8). Sometimes also He is said to be 'repentant,' inasmuch as in the eternal and immutable order of His providence, He builds up what He had previously destroyed, or destroys what He had previously made, as we do when moved by repentance: It repenteth me that I have made man (Gen. vi, 6, 7). God is also said to be 'sad,' inasmuch as things happen contrary to what He loves and approves, as sadness is in us at what happens against our will: And the Lord saw, and it seemed evil in his eyes, because judgement is not: God saw that there is no man, and he was displeased, because there was none to meet him (Isa. lix, 15, 16).

\section*{Chapter 92. In what sense Virtues can be posited in God}

\para[316] AS the divine goodness comprehends within itself in a certain way all goodnesses, and virtue is a sort of goodness, the divine goodness must contain all virtues after a manner proper to itself. But no virtue is predicated as an attribute of God after the manner of a habit, as virtues are in us. For it does not befit God to be good by anything superadded to Him, but only by His essence, since He is absolutely simple. Nor again does He act by anything superadded to His essence, as His essence is His being (Chap. ). Virtue therefore in God is not any habit, but His own essence.

\para[317] 2. A habit is an imperfect actuality, half-way between potentiality and actuality: hence the subjects of habits are compared to persons asleep. But in God actuality is most perfect. Virtue therefore in Him is not like a habit or a science, but is as a present act of consciousness, which is the extremest perfection of actuality.

\para[318] Since human virtues are for the guidance of human life, and human life is twofold, contemplative and active, the virtues of the active life, inasmuch as they perfect this present life, cannot be attributed to God: for the active life of man consists in the use of material goods, which are not assignable to God. Again, these virtues perfect human conduct in political society: hence they do not seem much to concern those who keep aloof from political society: much less can they befit God, whose conversation and life is far removed from the manner and custom of human life. Some again of the virtues of the active life direct us how to govern the passions: but in God there are no passions.

\section*{Chapter 93. That in God there are the Virtues which regulate Action}

\para[319] THERE are virtues directing the active life of man, which are not concerned with passions, but with actions, as truth, justice, liberality, magnificence, prudence, art. Since virtue is specified by its object, and the actions which are the objects of these virtues are not inconsistent with the divine perfection, neither is there in such virtues, specifically considered, anything to exclude them from the perfection of God.

\para[320] 3. Of things that come to have being from God, the proper plan of them all is in the divine understanding (Chap. ). But the plan of a thing to be made in the mind of the maker is Art: hence the Philosopher says that Art is "the right notion of things to be made." There is therefore properly Art in God, and therefore it is said: Wisdom, artificer of all, taught me (Wisd. vii, 21).

\para[321] 4. Again, the divine will, in things outside God, is determined by His knowledge (Chap. ). But knowledge directing the will to act is Prudence: because, according to the Philosopher, Prudence is "the right notion of things to be done." There is therefore Prudence in God; and hence it is said: With him is prudence (Job xii, 13).

\para[322] 5. From the fact of God wishing anything, He wishes the requisites of that thing. But the points requisite to the perfection of each several thing are due to that thing: there is therefore in God Justice, the function of which is to distribute to each his own. Hence it is said: The Lord is just, and hath loved justice (Ps. x, 8).

\para[323] 6. As shown above (Chapp. , ), the last end, for the sake of which God wills all things, in no way depends on the means to that end, neither in point of being nor in point of well-be ing. Hence God does not wish to communicate His goodness for any gain that may accrue to Himself thereby, but simply because the mere communication befits Him as the fountain of goodness. But to give, not from any advantage expected from the gift, but out of sheer goodness and the fitness of giving, is an act of Liberality. God therefore is in the highest degree liberal; and, as Avicenna says, He alone can properly be called liberal: for every other agent but Him is in the way of gaining something by his action and intends so to gain. This His liberality the Scripture declares, saying: As thou openest thy hand, all things shall be filled with goodness (Ps. ciii, 28) ; and, Who giveth to all abundantly, and reproacheth not (James i, 5).

\para[324] 7. All things that receive being from God, necessarily bear His likeness, in so far as they are, and are good, and have their proper archetypes in the divine understanding (Chap. ). But this belongs to the virtue of Truth, that every one should manifest himself in his deeds and words for such as he really is. There is therefore in God the virtue of Truth. Hence, God is true (Rom. iii, 4); and, All thy ways are truth (Ps. cxviii, 151).

\para[325] In point of exchange, the proper act of commutative justice, justice does not befit God, since He receives no advantage from any one; hence, Who hath first given to him, and recompense shall be made him? (Rom. xi, 35;) and, Who bath given to me beforehand, that I may repay him? (Job xli, 2.) Still, in a metaphorical sense, we are said to give things to God, inasmuch as He takes kindly what we have to offer Him. Commutative justice therefore does not befit God, but only distributive justice. To judge of things to be done, or to give a thing, or make a distribution, is not proper to man alone, but belongs to any and every intellectual being. Inasmuch therefore as the aforesaid actions are considered in their generality, they have their apt place even in divinity: for as man is the distributer of human goods, as of money or honour, so is God of all the goods of the universe. The aforesaid virtues therefore are of wider extension in God than in man: for as the justice of man is to a city or family, so is the justice of God to the entire universe: hence the divine virtues are said to be archetypes of ours. But other virtues, which do not properly become God, have no archetype in the divine nature, but only, as is the case with corporeal things generally, in the divine wisdom, which contains the proper notions of all things.

\section*{Chapter 94. That the Contemplative (Intellectual) Virtues are in God}

\para[326] IF Wisdom consists in the knowledge of the highest causes; and God chiefly knows Himself, and knows nothing except by knowing Himself, as the first cause of all (Chap. ), it is evident that Wisdom ought to be attributed to God in the first place. Hence it is said: He is wise of heart (Job ix, 4.); and, All wisdom is of the Lord God, and hath been with him alway (Ecclus i, 1). The Philosopher also says at the beginning of his Metaphysics that Wisdom is a divine possession, not a human.

\para[327] 2. If Knowledge (Science) is an acquaintance with a thing through its proper cause, and God knows the order of all causes and effects, and thereby the several proper causes of individual things (Chapp. , ), it is manifest that Knowledge (Science) is properly in God; hence God is the Lord of sciences (I Kings ii, 3)

\para[328] 3. If the immaterial cognition of things, attained without discussion, is Understanding (Intuition), God has such a cognition of all things (Chap. ); and therefore there is in Him Understanding. Hence, He hath counsel and understanding (Job xii, 13).

\section*{Chapter 95. That God cannot will Evil}

\para[329] EVERY act of God is an act of virtue, since Ills virtue is His essence (Chap. ).

\para[330] 2. The will cannot will evil except by some error coming to be in the reason, at least in the matter of the particular choice there and then made. For as the object of the will is good, apprehended as such, the will cannot tend to evil unless evil be somehow proposed to it as good; and that cannot be without error. But in the divine cognition there can be no error (Chap. ).

\para[331] 3. God is the sovereign good, admitting no intermixture of evil (Chap. ).

\para[332] 4. Evil cannot befall the will except by its being turned away from its end. But the divine will cannot be turned away from its end, being unable to will except by willing itself (Chap. ). It cannot therefore will evil; and thus free will in it is naturally established in good. This is the meaning of the texts: God is faithful and without iniquity (Deut. xxxii, 4); Thine eyes are clean, O Lord, and thou canst not look upon iniquity (Hab. i, 13).

\section*{Chapter 96. That God hates nothing}

\para[333] AS love is to good, so is hatred to evil; we wish good to them whom we love, and evil to them whom we hate. If then the will of God cannot be inclined to evil, as has been shown (Chap. ), it is impossible for Him to hate anything.

\para[334] 2. The will of God tends to things other than Himself inasmuch as, by willing and loving His own being and goodness, He wishes it to be diffused as far as is possible by communication of His likeness. This then is what God wills in beings other than Himself, that there be in them the likeness of His goodness. Therefore God wills the good of everything, and hates nothing.

\para[335] 4. What is found naturally in all active causes, must be found especially in the Prime Agent. But all agents in their own way love the effects which they themselves produce, as parents their children, poets their own poems, craftsmen their works. Much more therefore is God removed from hating anything, seeing that He is cause of all.

\para[336] Hence it is said: Thou lovest all things that are, and hatest nothing of the things that Thou hast made (Wisd. xi, 25).

\para[337] Some things however God is said, to hate figuratively (similitudinarie), and that in two ways. The first way is this, that God, in loving things and willing their good to be, wills their evil not to be: hence He is said to have hatred of evils, for the things we wish not to be we are said to hate. So it is said: Think no evil in your hearts every one of you against his friend, and love no lying oath: for all these are things that I hate, saith the Lord (Zach. viii, 17). But none of these things are effects of creation: they are not as subsistent things, to which hatred or love properly attaches. The other way is by God's wishing some greater good, which cannot be without the privation of a lesser good; and thus He is said to hate, whereas it is more properly love. Thus inasmuch as He wills the good of justice, or of the order of the universe, which cannot be without the punishment or perishing of some, He is said to hate those beings whose punishment or perishing He wills, according to the text, Esau I have hated (Malach. i, 3); and, Thou hatest all who work Iniquity, thou wilt destroy all who utter falsehood: the man of blood and deceit the Lord shall abominate (Ps. v, 7).

\section*{Chapter 97. That God is Living}

\para[338] IT has been shown that God is intelligent and willing: but to understand and will are functions of a living being only.

\para[339] 2. Life is attributed to beings inasmuch as they appear to move of themselves, and not to be moved by another. Therefore things that seem to move of themselves, the moving powers of which the vulgar do not perceive, are figuratively said to live, as we speak of the 'living' (running) water of a flowing stream, but not so of a cistern or stagnant pool; and we call 'quicksilver' that which seems to have a motion of its own. This is mere popular speech, for properly those things alone move of themselves, which do so by virtue of their composition of a moving force and matter moved, as things with souls; hence these alone are properly said to live: all other things are moved by some external force, a generating force, or a force removing an obstacle, or a force of impact. And because sensible activities are attended with movement, by a further step everything that determines itself to its own modes of activity, even though unattended with movement, is said to live; hence to understand and desire and feel are vital actions. But God, of all beings, is determined to activity by none other than Himself, as He is prime agent and first cause; to Him therefore, of all beings, does it belong to live.

\para[340] 3. The divine being contains the perfection of all being (Chap. ). But living is perfect being; hence animate things in the scale of being take precedence of inanimate. With God then to be is to live. This too is confirmed by authority of divine Scripture: I will raise to heaven my hand, and swear by my right hand, and say: I live for ever (Deut. xxxii, 40): My heart and my flesh) have rejoiced in the living God (Ps. lxxiii, 3).

\section*{Chapter 98. That God is His own Life}

\para[341] IN living things, to live is to be: for a living thing is said to be alive inasmuch as it has a soul; and by that soul, as by its own proper form, it has being: living in fact is nothing else than living being, arising out of a living form. But, in God, Himself is His own being (Chap. ): Himself therefore is His own life.

\para[342] 2. To understand is to live: but God is His own act of understanding (Chap. ).

\para[343] 3. If God is living, there must be life in Him. If then He is not His own life, there will be something in Him that is not Himself, and thus He will be compound, --- a rejected conclusion (Chap. ).

\para[344] And this is the text: I am life (John xiv, 6).

\section*{Chapter 99. That the Life of God is everlasting}

\para[345] IT is impossible for God to cease to live, since Himself He is His own life (Chap. ).

\para[346] 2. Everything that at one time is and at another time is not, has existence through some cause. But the divine life has no cause, as neither has the divine being. God is therefore not at one time living and at another not living, but always lives.

\para[347] 3. In every activity the agent remains, although sometimes the activity passes in succession: hence in motion the moving body remains the same in subject throughout the whole course of the motion, although not the same in our consideration. Where then the action is the agent himself, nothing there can pass in succession, but all must be together at once. But God's act of understanding and living is God Himself (Chapp. , ): therefore His life has no succession, but is all together at once, and everlasting.

\para[348] Hence it is said: This is the true God and life everlasting (i John V, 20).

\section*{Chapter 100. That God is Happy}

\para[349] HAPPINESS is the proper good of every intellectual nature. Since then God is an intellectual being, happiness will be His proper good. But God in regard of His proper good is not as a being that is still tending to a proper good not yet possessed: that is the way with a nature changeable and in potentiality; but God is in the position of a being that already possesses its proper good. Therefore He not only desires happiness, as we do, but is in the enjoyment of happiness.

\para[350] 2. The thing above all others desired or willed by an intellectual nature is the most perfect thing in that nature, and that is its happiness. But the most perfect thing in each is its most perfect activity: for power and habit are perfected by activity: hence the Philosopher says that happiness is a perfect activity. Now the perfection of activity depends on four conditions. First, on its kind, that it be immanent in the agent. I call an activity 'immanent in the agent,' when nothing else comes of it besides the act itself: such are the acts of seeing and hearing: such acts are perfections of the agents whose acts they are, and may have a finality of their own in so far as they are not directed to the production of anything else as an end. On the other hand, any activity from which there results something done besides itself, is a perfection of the thing done, not of the doer: it stands in the relation of a means to an end, and therefore cannot be the happiness of an intellectual nature. Secondly, on the principle of activity, that it be an activity of the highest power: hence our happiness lies not in any activity of sense, but in an activity of intellect, perfected by habit. Thirdly, on the object of activity; and therefore our happiness consists in understanding the highest object of understanding. Fourthly, on the form of activity, that the action be perfect, easy, and agreeable. But the activity of God fulfils all these conditions: since it is (1) activity in the order of understanding; and (2) His understanding is the highest of faculties, not needing any habit to perfect it; and (3) His understanding is bent upon Himself, the highest of intelligible objects; and (4) He understands perfectly, without any difficulty, and with all delight. He is therefore happy.

\para[351] 3. Boethius says that happiness is a state made perfect by a gathering of all good things. But such is the divine perfection, which includes all perfection in one single view (Chapp. , ).

\para[352] 4. He is happy, who is sufficient for himself and wants nothing. But God has no need of other things, seeing that His perfection depends on nothing external to Himself; and when He wills other things for Himself as for an end, it is not that He needs them, but only that this reference befits His goodness.

\para[353] 5. It is impossible for God to wish for anything impossible (Chap. ). Again it is impossible for anything to come in to Him which as yet He has not, seeing that He is nowise in potentiality (Chap. ). Therefore He cannot wish to have what He has not: therefore He has whatever He wishes; and He wishes nothing evil (Chap. ). Therefore He is happy, according to the definition given by some, that "he is happy who has what he wishes and wishes nothing evil." His happiness the Holy Scriptures declare: Whom he will show in his own time, the blessed and powerful one (i Tim. vi, 15).

\section*{Chapter 101. That God Is His own Happiness}

\para[354] GOD'S happiness is the act of His understanding (Chap. ). But that very act of God's understanding is His substance (Chap. ). He therefore is His own happiness.

\section*{Chapter 102. That the Happiness of God is most perfect, and exceeds all other happiness}

\para[355] WHERE there is greater love, there is greater delight in the attainment of the object loved. But every being, other things being equal, loves itself more than it loves anything else: a sign of which is that, the nearer anything is to oneself, the more it is naturally loved. God therefore takes greater delight in His happiness, which is Himself, than other blessed ones in their happiness, which is not what they are.

\para[356] 3. What is by essence, ranks above what is by participation. But God is happy by His essence, a prerogative that can belong to no other: for nothing else but God can be the sovereign good; and thus whatever else is happy must be happy by participation from Him. The divine happiness therefore exceeds all other happiness.

\para[357] 4. Perfect happiness consists in an act of the understanding. But no other act of understanding can compare with God's act: as is clear, not only from this that it is a subsistent act, but also because by this one act God perfectly understands Himself as He is, and all things that are and are not, good and evil; whereas in all other intellectual beings the act of understanding is not itself subsistent, but is the act of a subsistent subject. Nor can any one understand God, the supreme object of understanding, so perfectly as He is perfect, because the being of none is so perfect as the divine being, nor can any act ever be more perfect than the substance of which it is the act. Nor is there any other understanding that knows even all that God can do: for if it did, it would comprehend the divine power. Lastly, even what another understanding does know, it does not know all with one and the same act. God therefore is incomparably happy above all other beings.

\para[358] The more a thing is brought to unity, the more perfect is its power and excellence. But an activity that works in succession, is divided by different divisions of time: in no way then can its perfection be compared to the perfection of an activity that is without succession, all present together, especially if it does not pass in an instant but abides to eternity. Now the divine act of understanding is without succession, existing all together for eternity: whereas our act of understanding is in succession by the accidental attachment to it of continuity and time. Therefore the divine happiness infinitely exceeds human happiness, as the duration of eternity exceeds the 'now in flux' of time (nunc temporis fluens).

\para[359] 6. The fatigue and various occupations whereby our contemplation in this life is necessarily interrupted, --- in which contemplation whatever happiness there is for man in this life chiefly consists, --- and the errors and doubts and various mishaps to which the present life is subject, show that human happiness, in this life particularly, can in no way compare with the happiness of God.

\para[360] 7. The perfection of the divine happiness may be gathered from this, that it embraces all happinesses according to the most perfect mode of each. By way of contemplative happiness, it has a perfect and perpetual view of God Himself and of other beings. By way of active life, it has the government, not of one man, or of one house, or of one city, or of one kingdom, but of the whole universe. Truly, the false happiness of earth is but a shadow of that perfect happiness. For it consists, according to Boethius, in five things, in pleasure, riches, power, dignity and fame. God then has a most excellent delight of Himself, and a universal joy of all good things, without admixture of contrary element. For riches, He has absolute self-sufficiency of all good. For power, He has infinite might. For dignity, He has primacy and rule over all beings. For fame, He has the admiration of every understanding that in any sort knows Him. To Him then, who is singularly blessed, be honour and glory for ever and ever, Amen.
% !body-end
\end{document}
